% ============================================================================
% PART I: MATHEMATICAL FOUNDATIONS (fragment; preamble lives in master file)
% ============================================================================
\part{Mathematical Foundations}

Before describing the \nm{} mechanism itself, we assemble the mathematical
toolkit its theory rests on. Four ingredients recur throughout. First,
\emph{decision-focused learning}: \nm{} never scores a forecast by how close it
is to the truth, but by the portfolio decision it induces, so we need the exact
geometry of the decision layer, the induced regret, and a trainable surrogate
(SPO+). Second, \emph{concentration of measure}: every retention and transfer
claim ultimately reduces to ``a quantity estimated from $5$--$7$ episodes is
close to its mean,'' and we must be honest about what such small samples can
and cannot certify. Third, \emph{distances between distributions}: the
regime-conditional guarantees we import are stated in divergence
(KL/total-variation) terms, and a small but load-bearing ``Pinsker bridge''
converts them into statements about bounded decision regret. Fourth,
\emph{online learning on the simplex}: the gate that routes episodes to niches
is an exponentiated-gradient update, and we must understand both what its
classical regret theory gives us and---just as importantly---what it does
\emph{not} give us once losses become endogenous. The final section formalizes
why pooled empirical risk minimization fails across non-stationary episodes and
why an across-episode variance penalty detects exactly the failure mode,
positioning the \nm{} objective inside the invariant-risk literature.

% ============================================================================
\section{Decision-Focused Learning from First Principles}
\label{sec:dd1-dfl}

\subsection{Prediction Versus Decision: the Predict-then-Optimize Pipeline}
\label{sec:dd1-pto}

Classical supervised learning evaluates a forecast $\hat{y}$ of an unknown
quantity $y$ by a \emph{prediction loss} such as the squared error
$\|\hat{y}-y\|_2^2$. But in \nm{}---as in a large class of operational
problems---the forecast is not the product. The forecast is an
\emph{intermediate input} to an optimization problem whose solution is the
actual decision, and the decision is what earns or loses money. This is the
\textbf{predict-then-optimize} pipeline \cite{elmachtoub2022smart,
mandi2020smart, mandi2024decision}:
\begin{equation}
\label{eq:dd1-pipeline}
x \;\xrightarrow{\ \text{model}\ }\; \hat{y} = f_\theta(x)
\;\xrightarrow{\ \text{optimizer}\ }\; z^*(\hat{y}) \in \argmax_{z \in Z}
F(z, \hat{y})
\;\xrightarrow{\ \text{nature reveals } y\ }\; \text{realized value } F(z^*(\hat{y}), y).
\end{equation}
Here $x$ are observable features (market state), $y \in \R^n$ is the unknown
parameter vector (next-period asset returns), $Z \subseteq \R^n$ is a known
feasible set of decisions (portfolios), and $F : Z \times \R^n \to \R$ is the
objective.

\begin{definition}[Linear-objective decision layer]
\label{def:dd1-linear-layer}
A decision layer is \textbf{linear-objective} if $F(z, y) = y^\top z$, i.e.\
the unknown parameter enters the objective linearly, while all problem
structure (combinatorics, budgets, bounds) is carried by the feasible set $Z$,
which does not depend on $y$. We write
\begin{equation}
\label{eq:dd1-oracle-def}
z^*(y) \in \argmax_{z \in Z} \, y^\top z, \qquad
V(y) := \max_{z \in Z} \, y^\top z ,
\end{equation}
and call the map $y \mapsto z^*(y)$ the \textbf{decision oracle} and $V$ the
\textbf{support function} of $Z$.
\end{definition}

\begin{remark}
Linearity of the objective in $y$ is a genuine restriction, but it covers
shortest paths, assignment, knapsack relaxations, and---our case---linear
portfolio construction. It buys us two things: the optimal value $V(y)$ is a
maximum of linear functions of $y$ and hence \emph{convex} in $y$, and the SPO+
surrogate of Section~\ref{sec:dd1-spoplus} exists in closed form.
\end{remark}

\nm{}'s decision layer is a cardinality-constrained, long-only allocation.

\begin{definition}[Cardinality-constrained long-only portfolio]
\label{def:dd1-portfolio}
Fix the universe size $n \in \N$, a per-asset cap $w_{\max} > 0$, a gross
budget $W > 0$, and a cardinality $k \in \{1, \dots, n\}$. The feasible set is
\begin{equation}
\label{eq:dd1-Z}
Z \;=\; \Bigl\{\, z \in [0, w_{\max}]^n \;:\; \textstyle\sum_{i=1}^n z_i \le W,
\;\; \|z\|_0 \le k \,\Bigr\},
\end{equation}
where $\|z\|_0 = |\{i : z_i \neq 0\}|$ counts nonzero positions. The objective
is the portfolio return $F(z, y) = y^\top z$ for a return vector $y \in \R^n$.
\end{definition}

Note that $Z$ is \emph{nonconvex} (the $\ell_0$ constraint carves the box into
$\binom{n}{k}$-many faces), so generic convex-optimization machinery does not
apply. Fortunately the problem admits an exact, closed-form, sort-based
solution.

\begin{theorem}[Exact decision oracle]
\label{thm:dd1-oracle}
Assume $W \ge k\, w_{\max}$ (the budget never binds before the cap and the
cardinality do). For $y \in \R^n$ let
\[
P(y) = \{ i : y_i > 0 \}, \qquad
S^*(y) = \text{the } \min\{k, |P(y)|\} \text{ indices of } P(y) \text{ with
largest } y_i,
\]
with ties broken toward the smaller index. Then
\begin{equation}
\label{eq:dd1-oracle-form}
z^*(y)_i \;=\; w_{\max} \,\mathbf{1}\{ i \in S^*(y) \}
\end{equation}
is an optimal solution of $\max_{z \in Z} y^\top z$, with optimal value
$V(y) = w_{\max} \sum_{i \in S^*(y)} y_i \ge 0$.
\end{theorem}

\begin{proof}
\emph{Feasibility.} The vector in \eqref{eq:dd1-oracle-form} has entries in
$\{0, w_{\max}\} \subseteq [0, w_{\max}]$, support size $|S^*(y)| \le k$, and
gross exposure $|S^*(y)|\, w_{\max} \le k\, w_{\max} \le W$. So
$z^*(y) \in Z$.

\emph{Optimality.} Let $z \in Z$ be arbitrary with support
$T = \{i : z_i > 0\}$, so $|T| \le k$. Using $0 \le z_i \le w_{\max}$ and
writing $(u)_+ = \max\{u, 0\}$,
\[
y^\top z \;=\; \sum_{i \in T} y_i z_i
\;\le\; \sum_{i \in T} (y_i)_+ \, z_i
\;\le\; w_{\max} \sum_{i \in T} (y_i)_+ .
\]
The first inequality holds termwise: if $y_i \le 0$ then $y_i z_i \le 0 =
(y_i)_+ z_i$; if $y_i > 0$ the two sides agree. Now
$\sum_{i \in T} (y_i)_+$ is a sum of at most $k$ of the nonnegative numbers
$\{(y_i)_+\}_{i=1}^n$, hence is at most the sum of the $k$ largest of them,
which equals $\sum_{i \in S^*(y)} y_i$ (indices outside $P(y)$ contribute
$(y_i)_+ = 0$ and are never needed). Therefore
$y^\top z \le w_{\max} \sum_{i \in S^*(y)} y_i = y^\top z^*(y)$. Finally
$V(y) \ge 0$ because $z = 0 \in Z$.
\end{proof}

\begin{keypoint}
The oracle is a \textbf{sort}: rank the forecast coordinates, keep the top $k$
strictly positive ones, allocate the cap $w_{\max}$ to each. It costs
$O(n \log n)$, is exact (no relaxation gap despite the nonconvex $\ell_0$
constraint), and will be called \emph{twice per gradient step} by the SPO+
subgradient of Section~\ref{sec:dd1-spoplus}. The standing assumption
$W \ge k\, w_{\max}$ is maintained throughout the document; when it fails the
oracle becomes a fractional-knapsack rule and only the boundary case changes.
\end{keypoint}

Throughout, we fix the tie-breaking rule above so that $z^*(\cdot)$ is a
well-defined \emph{function} (not a set-valued map); all statements below are
for this selection.

\subsection{Decision Regret and Its Geometry}
\label{sec:dd1-regret}

\begin{definition}[Decision regret]
\label{def:dd1-regret}
For a forecast $\hat{y} \in \R^n$ and realized returns $y \in \R^n$, the
\textbf{decision regret} (also: SPO loss \cite{elmachtoub2022smart}) is
\begin{equation}
\label{eq:dd1-regret}
\rho(\hat{y}; y) \;=\; F\bigl(z^*(y), y\bigr) - F\bigl(z^*(\hat{y}), y\bigr)
\;=\; y^\top z^*(y) - y^\top z^*(\hat{y}).
\end{equation}
It is the return given up by acting on the forecast instead of acting on the
truth.
\end{definition}

\begin{proposition}[Range of the regret]
\label{prop:dd1-regret-range}
Suppose $\|y\|_\infty \le y_{\max}$. Then for every $\hat{y} \in \R^n$,
\begin{equation}
\label{eq:dd1-B}
0 \;\le\; \rho(\hat{y}; y) \;\le\; B, \qquad B := 2\, k\, w_{\max}\, y_{\max}.
\end{equation}
\end{proposition}

\begin{proof}
\emph{Lower bound.} $z^*(\hat{y}) \in Z$, and $z^*(y)$ maximizes $y^\top z$
over $Z$ (Theorem~\ref{thm:dd1-oracle}), so
$y^\top z^*(y) \ge y^\top z^*(\hat{y})$.

\emph{Upper bound.} For any $z \in Z$ with support $T$, $|T| \le k$ and
$0 \le z_i \le w_{\max}$, so
\[
\bigl| y^\top z \bigr| \;\le\; \sum_{i \in T} |y_i|\, z_i
\;\le\; k\, w_{\max}\, y_{\max}.
\]
Hence $y^\top z^*(y) \le k\, w_{\max}\, y_{\max}$ and
$y^\top z^*(\hat{y}) \ge -\,k\, w_{\max}\, y_{\max}$; subtracting gives
$\rho \le 2 k\, w_{\max}\, y_{\max}$.
\end{proof}

\begin{remark}
The constant $B$ is attained only in a worst case where the forecast selects
$k$ assets that each lose $y_{\max}$ while the best portfolio gains $y_{\max}$
on each of $k$ assets. With daily equity returns this is astronomically
pessimistic; we return to this honestly in
Remark~\ref{rem:dd1-B-loose} when $B$ enters the Pinsker bridge.
\end{remark}

The next observation is small but conceptually central: the decision layer
destroys all information about the \emph{scale} of the forecast.

\begin{proposition}[Positive scale invariance; regret is a ranking loss]
\label{prop:dd1-scale-invariance}
For every $c > 0$ and $\hat{y} \in \R^n$:
\begin{enumerate}
\item $z^*(c\,\hat{y}) = z^*(\hat{y})$, and consequently
$\rho(c\,\hat{y}; y) = \rho(\hat{y}; y)$.
\item More generally, $\rho(\hat{y}; y)$ depends on $\hat{y}$ only through
the sign pattern $(\mathrm{sign}(\hat{y}_i))_{i}$ and the relative order of
the positive coordinates of $\hat{y}$.
\end{enumerate}
\end{proposition}

\begin{proof}
(1) Multiplication by $c > 0$ preserves the positivity set,
$P(c\hat{y}) = P(\hat{y})$, and preserves all pairwise comparisons
$c\hat{y}_i > c\hat{y}_j \iff \hat{y}_i > \hat{y}_j$ (including ties, which
are broken by the same index rule). Hence $S^*(c\hat{y}) = S^*(\hat{y})$ and,
by \eqref{eq:dd1-oracle-form}, $z^*(c\hat{y}) = z^*(\hat{y})$. The regret
\eqref{eq:dd1-regret} depends on $\hat{y}$ only through $z^*(\hat{y})$.

(2) Inspecting Theorem~\ref{thm:dd1-oracle}: the selected set $S^*(\hat{y})$
is determined by which coordinates are strictly positive and by the ordering
among those, and nothing else. Since $\rho(\cdot\,; y)$ factors through
$S^*(\hat{y})$, the claim follows.
\end{proof}

\begin{intuition}
The optimizer asks the forecast exactly one question: \emph{``which $k$ assets
do you like, and do you like them at all?''} Calibration, magnitude, even the
gap between the best and second-best asset are invisible to the decision. A
loss that trains the forecaster should therefore reward correct
\emph{rankings}, not small residuals---this is the entire raison d'\^etre of
decision-focused learning \cite{elmachtoub2022smart, shah2022lodl}.
\end{intuition}

\begin{example}[Zero MSE correlation with regret: a 4-asset case]
\label{ex:dd1-mse-vs-regret}
Let $n = 4$, $k = 2$, $w_{\max} = 0.5$, $W = 1$, and realized returns
\[
y = (0.05,\; 0.03,\; -0.02,\; 0.01).
\]
The oracle selects the two largest positive coordinates: $S^*(y) = \{1, 2\}$,
$z^*(y) = (0.5, 0.5, 0, 0)$, and the optimal value is
$V(y) = 0.5\,(0.05 + 0.03) = 0.04$.

\textbf{Forecast A (high MSE, zero regret).} Take $\hat{y}^A = 10\,y =
(0.5, 0.3, -0.2, 0.1)$, a wildly miscalibrated forecast. Its squared error is
\[
\mathrm{MSE}(\hat{y}^A)
= \tfrac{1}{4}\bigl(0.45^2 + 0.27^2 + 0.18^2 + 0.09^2\bigr)
= \tfrac{1}{4}(0.2025 + 0.0729 + 0.0324 + 0.0081)
\approx 0.0790 .
\]
But by Proposition~\ref{prop:dd1-scale-invariance},
$z^*(\hat{y}^A) = z^*(y)$, so $\rho(\hat{y}^A; y) = 0$: the decision is
\emph{perfect}.

\textbf{Forecast B (tiny MSE, large regret).} Take
$\hat{y}^B = (0.025,\; 0.030,\; -0.020,\; 0.031)$. Its squared error is
\[
\mathrm{MSE}(\hat{y}^B)
= \tfrac{1}{4}\bigl(0.025^2 + 0^2 + 0^2 + 0.021^2\bigr)
= \tfrac{1}{4}(0.000625 + 0.000441)
\approx 0.000267,
\]
about $296$ times smaller than forecast A's. Yet the induced ranking puts
assets $4$ and $2$ on top: $S^*(\hat{y}^B) = \{4, 2\}$,
$z^*(\hat{y}^B) = (0, 0.5, 0, 0.5)$, realized value
$y^\top z^*(\hat{y}^B) = 0.5\,(0.03 + 0.01) = 0.02$, and
\[
\rho(\hat{y}^B; y) = 0.04 - 0.02 = 0.02,
\]
i.e.\ \emph{half} the achievable return is forfeited. MSE ranks B far above
A; the decision ranks A strictly above B.
\end{example}

\begin{warning}
Example~\ref{ex:dd1-mse-vs-regret} is not a pathological corner case. In a
cross-section of assets the return differences that determine the top-$k$ set
are routinely smaller than the irreducible noise level, so an MSE-optimal
forecast spends its capacity matching magnitudes that the decision layer
ignores, while leaving the decision-relevant orderings to chance. This is the
failure that the SPO family of losses is designed to repair.
\end{warning}

\subsection{The SPO Loss, Its Pathologies, and the SPO+ Surrogate}
\label{sec:dd1-spoplus}

The natural training loss is the regret itself, viewed as a function of the
prediction: $\ell_{\mathrm{SPO}}(\hat{y}; y) := \rho(\hat{y}; y)$. It is,
unfortunately, almost unusable directly.

\begin{example}[Discontinuity and nonconvexity of the SPO loss]
\label{ex:dd1-spo-discont}
Let $n = 2$, $k = 1$, $w_{\max} = 1$, $W \ge 1$, $y = (0.04, 0.02)$, and
consider the forecast family $\hat{y}(t) = (t, 0.02)$ for $t \in \R$. The
oracle picks asset $1$ iff $t > 0.02$ (at $t = 0.02$ the tie-break also picks
asset $1$), and asset $2$ for $0 < t < 0.02$. Hence
\[
\ell_{\mathrm{SPO}}\bigl(\hat{y}(t); y\bigr) =
\begin{cases}
0, & t \ge 0.02 \quad (\text{decision } z = e_1,\ \text{value } 0.04),\\[2pt]
0.02, & t < 0.02 \quad (\text{decision } z = e_2,\ \text{value } 0.02),
\end{cases}
\]
a step function: \textbf{discontinuous} at $t = 0.02$, with derivative $0$
everywhere else. Any nonconstant function that is piecewise constant is
nonconvex (a finite convex function on $\R$ is continuous, and a continuous
piecewise-constant function is constant), so $\ell_{\mathrm{SPO}}$ is also
\textbf{nonconvex} in $\hat{y}$. Gradient-based training receives either no
signal (zero gradient) or an undefined one (at the jump).
\end{example}

The remedy of \cite{elmachtoub2022smart} is a convex surrogate, SPO+. Their
derivation is stated for cost \emph{minimization}; \nm{} maximizes return, so
we derive the maximization form explicitly via the substitution $c = -y$.

\paragraph{Derivation.} In the minimization convention one has cost vectors
$c$, decisions $w^*(c) \in \argmin_{w \in Z} c^\top w$, and the SPO+ loss
\begin{equation}
\label{eq:dd1-spoplus-min}
\ell_{\mathrm{SPO+}}^{\min}(\hat{c}; c)
\;=\; \max_{w \in Z}\,\bigl\{ (c - 2\hat{c})^\top w \bigr\}
\;+\; 2\,\hat{c}^\top w^*(c) \;-\; c^\top w^*(c),
\end{equation}
obtained from the SPO loss by a duality relaxation and the (provably
tightest) scaling $\hat{c} \mapsto 2\hat{c}$ \cite{elmachtoub2022smart}.
Substituting $c = -y$ and $\hat{c} = -\hat{y}$, and noting that
$\argmin_{w} (-y)^\top w = \argmax_{z} y^\top z$ so $w^*(c) = z^*(y)$, the
three terms of \eqref{eq:dd1-spoplus-min} become
$\max_z (2\hat{y} - y)^\top z$, $-2\hat{y}^\top z^*(y)$, and
$+\,y^\top z^*(y)$. We adopt the result as our definition.

\begin{definition}[SPO+ loss, maximization form]
\label{def:dd1-spoplus}
For forecast $\hat{y}$ and realized returns $y$,
\begin{equation}
\label{eq:dd1-spoplus}
\ell_{\mathrm{SPO+}}(\hat{y}; y)
\;=\; \underbrace{\max_{z \in Z}\, (2\hat{y} - y)^\top z}_{\textstyle
= \,V(2\hat{y} - y)}
\;-\; 2\,\hat{y}^\top z^*(y)
\;+\; y^\top z^*(y).
\end{equation}
Both extremal problems are instances of the same oracle
(Theorem~\ref{thm:dd1-oracle}), applied to the \emph{perturbed} vector
$2\hat{y} - y$ and to the true vector $y$.
\end{definition}

A reassuring sanity check: at the truth, the surrogate vanishes,
\[
\ell_{\mathrm{SPO+}}(y; y) = V(y) - 2\,y^\top z^*(y) + y^\top z^*(y)
= V(y) - y^\top z^*(y) = 0 .
\]

\begin{proposition}[Convexity in the prediction]
\label{prop:dd1-spoplus-convex}
For every fixed $y$, the map $\hat{y} \mapsto \ell_{\mathrm{SPO+}}(\hat{y}; y)$
is convex on $\R^n$ (indeed, piecewise affine).
\end{proposition}

\begin{proof}
The first term is $\hat{y} \mapsto V(2\hat{y} - y) =
\max_{z \in Z} \{ 2 z^\top \hat{y} - z^\top y \}$, a pointwise maximum of
affine functions of $\hat{y}$, hence convex (a supremum of convex functions is
convex, and affine functions are convex). The second term
$-2\,\hat{y}^\top z^*(y)$ is affine in $\hat{y}$ ($z^*(y)$ is a constant once
$y$ is fixed), and the third term is a constant. A sum of convex functions is
convex. Since $Z$ has finitely many extreme points (each coordinate of an
optimal $z$ is $0$ or $w_{\max}$ on one of finitely many supports), the
maximum is over finitely many affine functions, so the loss is piecewise
affine.
\end{proof}

\begin{proposition}[Majorization: SPO+ upper-bounds the decision regret]
\label{prop:dd1-spoplus-majorize}
For all $\hat{y}, y \in \R^n$,
\begin{equation}
\label{eq:dd1-majorize}
\rho(\hat{y}; y) \;\le\; \ell_{\mathrm{SPO+}}(\hat{y}; y).
\end{equation}
\end{proposition}

\begin{proof}
Lower-bound the maximum in \eqref{eq:dd1-spoplus} by the particular feasible
point $z = z^*(\hat{y})$:
\[
\ell_{\mathrm{SPO+}}(\hat{y}; y)
\;\ge\; (2\hat{y} - y)^\top z^*(\hat{y}) - 2\,\hat{y}^\top z^*(y)
+ y^\top z^*(y)
\;=\; 2\,\hat{y}^\top z^*(\hat{y}) - y^\top z^*(\hat{y})
- 2\,\hat{y}^\top z^*(y) + y^\top z^*(y).
\]
Now $z^*(\hat{y})$ maximizes $\hat{y}^\top z$ over $Z$, so
$\hat{y}^\top z^*(\hat{y}) \ge \hat{y}^\top z^*(y)$, i.e.\
$2\,\hat{y}^\top z^*(\hat{y}) - 2\,\hat{y}^\top z^*(y) \ge 0$. Dropping this
nonnegative quantity,
\[
\ell_{\mathrm{SPO+}}(\hat{y}; y)
\;\ge\; y^\top z^*(y) - y^\top z^*(\hat{y})
\;=\; \rho(\hat{y}; y). \qedhere
\]
\end{proof}

Combining Propositions~\ref{prop:dd1-regret-range} and
\ref{prop:dd1-spoplus-majorize}: minimizing the convex surrogate pushes down
the true (discontinuous) regret. The surrogate is moreover statistically
faithful in the following sense.

\begin{theorem}[Fisher consistency of SPO+; \cite{elmachtoub2022smart}, stated
without proof]
\label{thm:dd1-fisher}
Suppose the conditional distribution of $y$ given $x$ is continuous, centrally
symmetric about its mean $\bar{y}(x) = \E[y \mid x]$, and assigns the optimal
decision uniquely for almost every realization. Then the minimizer of the
conditional SPO+ risk $\E\bigl[\ell_{\mathrm{SPO+}}(\hat{y}; y) \mid x\bigr]$
over $\hat{y} \in \R^n$ is $\hat{y} = \bar{y}(x)$; in particular, minimizing
SPO+ risk also minimizes the true SPO (regret) risk.
\end{theorem}

\paragraph{Subgradient.} Convexity (Proposition~\ref{prop:dd1-spoplus-convex})
gives the loss a nonempty subdifferential everywhere; the support-function
structure makes one element of it available in closed form. We first record
the elementary lemma usually attributed to Danskin in this finite setting.

\begin{lemma}[Subgradient of a support function]
\label{lem:dd1-danskin}
Let $h(u) = \max_{z \in Z} u^\top z$ with $Z$ compact, and let
$\bar{z} \in \argmax_{z \in Z} u^\top z$. Then $\bar{z} \in \partial h(u)$,
i.e.\ $h(u') \ge h(u) + \bar{z}^\top (u' - u)$ for all $u'$.
\end{lemma}

\begin{proof}
$h(u') \ge u'^\top \bar{z} = u^\top \bar{z} + (u' - u)^\top \bar{z}
= h(u) + \bar{z}^\top (u' - u)$, where the first step holds because $\bar{z}$
is feasible for the maximization defining $h(u')$ and the last because
$\bar{z}$ attains $h(u)$.
\end{proof}

\begin{proposition}[SPO+ subgradient and the linear-model chain rule]
\label{prop:dd1-spoplus-grad}
Fix $y$ and let $\hat{y} \in \R^n$. Then
\begin{equation}
\label{eq:dd1-subgrad}
g(\hat{y}; y) \;:=\; 2\,\Bigl( z^*\bigl(2\hat{y} - y\bigr) - z^*(y) \Bigr)
\;\in\; \partial_{\hat{y}}\, \ell_{\mathrm{SPO+}}(\hat{y}; y),
\end{equation}
so a subgradient-descent step moves the forecast in the direction
\begin{equation}
\label{eq:dd1-descent-dir}
-\,g(\hat{y}; y) \;=\; 2\,\Bigl( z^*(y) - z^*\bigl(2\hat{y} - y\bigr) \Bigr),
\end{equation}
i.e.\ \emph{toward} the assets the truth selects and \emph{away from} the
assets the perturbed forecast selects. If the forecast is produced by a linear
model $\hat{y} = \Theta x$ with $\Theta \in \R^{n \times d}$ and features
$x \in \R^d$, then by the affine chain rule for subgradients,
\begin{equation}
\label{eq:dd1-chain}
g(\Theta x; y)\, x^\top \;\in\; \partial_{\Theta}\,
\ell_{\mathrm{SPO+}}(\Theta x;\, y) \;\subseteq\; \R^{n \times d},
\end{equation}
and for a batch $\{(x_s, y_s)\}_{s=1}^m$ the subgradient of the average loss
is $\frac{1}{m} \sum_s g(\Theta x_s; y_s)\, x_s^\top$.
\end{proposition}

\begin{proof}
Write $\ell_{\mathrm{SPO+}}(\hat{y}; y) = h(2\hat{y} - y) - 2\,\hat{y}^\top
z^*(y) + \text{const}$, with $h$ the support function of $Z$. By
Lemma~\ref{lem:dd1-danskin}, $z^*(2\hat{y} - y) \in \partial h(2\hat{y} - y)$
(our oracle applied to the vector $2\hat{y} - y$ is exactly an argmax for
$h$). The map $\hat{y} \mapsto 2\hat{y} - y$ is affine with Jacobian $2 I$, so
by the affine chain rule $2\,z^*(2\hat{y} - y)$ is a subgradient of
$\hat{y} \mapsto h(2\hat{y} - y)$. The term $-2\,\hat{y}^\top z^*(y)$ is
differentiable with gradient $-2\,z^*(y)$. Subdifferentials of sums of convex
functions (one of them differentiable) add, giving \eqref{eq:dd1-subgrad}.
The chain rule \eqref{eq:dd1-chain} is the same affine-composition rule
applied to $\Theta \mapsto \Theta x$: for any $\Theta'$,
$\ell(\Theta' x) \ge \ell(\Theta x) + g^\top(\Theta' - \Theta)x
= \ell(\Theta x) + \langle g x^\top, \Theta' - \Theta\rangle_F$.
\end{proof}

\begin{remark}[Sign conventions]
\label{rem:dd1-sign}
In the cost-minimization convention of \cite{elmachtoub2022smart} the same
object reads $2\bigl(w^*(c) - w^*(2\hat{c} - c)\bigr)$. Under $c = -y$ the two
expressions map onto each other; the expression
$2\bigl(z^*(y) - z^*(2\hat{y} - y)\bigr)$ that appears in informal
descriptions of \nm{} is the \emph{descent direction}
\eqref{eq:dd1-descent-dir} of the maximization form, not the subgradient
itself. Both conventions produce the identical parameter update; only the
bookkeeping differs.
\end{remark}

\begin{intuition}
The update \eqref{eq:dd1-descent-dir} has a ``contrastive'' reading. The
oracle is run on the truth $y$ and on the reflection $2\hat{y} - y$ of the
truth through the forecast (the point such that $\hat{y}$ is the midpoint of
$y$ and $2\hat{y} - y$; the reflection exaggerates the forecast's deviation
from the truth). Assets selected by the truth but not by the exaggerated
forecast get their predicted returns pushed \emph{up} by $2 w_{\max} \eta$ per
step; assets selected only by the exaggerated forecast get pushed \emph{down}
by the same amount; assets on which both agree (or both pass) receive no
gradient. The loss is exactly zero gradient when the exaggerated forecast
already makes the right selection---training stops adjusting magnitudes the
moment the ranking is robustly correct, which is precisely the
scale-invariance of Proposition~\ref{prop:dd1-scale-invariance} made
operational.
\end{intuition}

\begin{example}[One SPO+ gradient step on four assets]
\label{ex:dd1-spoplus-step}
Continue with $n = 4$, $k = 2$, $w_{\max} = 0.5$, $W = 1$, true returns
$y = (0.05, 0.03, -0.02, 0.01)$, and current forecast
\[
\hat{y} = (0.02,\; 0.01,\; 0.04,\; 0.00).
\]

\emph{Step 1: the two oracle calls.} The truth selects
$S^*(y) = \{1, 2\}$, so $z^*(y) = (0.5, 0.5, 0, 0)$. The perturbed vector is
\[
2\hat{y} - y = (0.04 - 0.05,\; 0.02 - 0.03,\; 0.08 + 0.02,\; 0 - 0.01)
= (-0.01,\; -0.01,\; 0.10,\; -0.01),
\]
whose only positive coordinate is the third, so
$S^*(2\hat{y} - y) = \{3\}$ and $z^*(2\hat{y} - y) = (0, 0, 0.5, 0)$.

\emph{Step 2: loss value.} By \eqref{eq:dd1-spoplus},
\[
\ell_{\mathrm{SPO+}} = \underbrace{0.10 \times 0.5}_{V(2\hat y - y) = 0.05}
\;-\; \underbrace{2\,(0.02 + 0.01)(0.5)}_{2\hat{y}^\top z^*(y) = 0.03}
\;+\; \underbrace{0.04}_{y^\top z^*(y)} \;=\; 0.06 .
\]
For comparison, the true regret of $\hat{y}$: the forecast selects its top
two positive coordinates $S^*(\hat{y}) = \{3, 1\}$, realized value
$0.5\,(-0.02 + 0.05) = 0.015$, so
$\rho(\hat{y}; y) = 0.04 - 0.015 = 0.025 \le 0.06$, consistent with
Proposition~\ref{prop:dd1-spoplus-majorize}.

\emph{Step 3: subgradient and update.} By \eqref{eq:dd1-subgrad},
\[
g = 2\bigl( (0, 0, 0.5, 0) - (0.5, 0.5, 0, 0) \bigr)
= (-1,\; -1,\; 1,\; 0).
\]
A descent step with rate $\eta = 0.02$ gives
\[
\hat{y}' = \hat{y} - \eta\, g
= (0.02 + 0.02,\; 0.01 + 0.02,\; 0.04 - 0.02,\; 0.00)
= (0.04,\; 0.03,\; 0.02,\; 0.00).
\]
The new forecast ranks the assets $1 \succ 2 \succ 3 \succ 4$, selects
$S^*(\hat{y}') = \{1, 2\}$, and achieves $\rho(\hat{y}'; y) = 0$: one step
flipped the decision from losing $0.025$ to optimal. Note what the step did
\emph{not} do: it made the forecast magnitudes \emph{less} accurate for asset
$3$ in MSE terms en route to fixing the ranking.
\end{example}

\begin{keypoint}
Summary of the decision-focused toolkit used by \nm{}:
\begin{itemize}
\item the decision layer is an $O(n \log n)$ sort oracle
      (Theorem~\ref{thm:dd1-oracle});
\item the quantity that matters is the regret
      $\rho \in [0, B]$, a bounded \emph{ranking} loss
      (Propositions~\ref{prop:dd1-regret-range},
      \ref{prop:dd1-scale-invariance});
\item $\rho$ itself is untrainable
      (Example~\ref{ex:dd1-spo-discont}), so each niche head is trained on
      the convex majorant $\ell_{\mathrm{SPO+}}$
      (Definition~\ref{def:dd1-spoplus}), whose subgradient costs two oracle
      calls (Proposition~\ref{prop:dd1-spoplus-grad}).
\end{itemize}
Every ``episode risk'' appearing later in the document is an average of
$\ell_{\mathrm{SPO+}}$ (for training) or of $\rho$ (for evaluation) over the
days of an episode.
\end{keypoint}

% ============================================================================
\section{Concentration of Measure}
\label{sec:dd1-concentration}

\nm{}'s retention and transfer arguments repeatedly take the form: ``the
average regret of niche head $r$ over its $E_r$ stored episodes is small;
therefore its regret on a fresh episode of the same regime is small with high
probability.'' Since $E_r$ is tiny ($5$--$7$ in the companion experiments),
we need the sharpest available finite-sample inequalities---and an honest
accounting of how weak even those are at such sample sizes.

\subsection{Markov, Chebyshev, and Cantelli}

\begin{theorem}[Markov's inequality]
\label{thm:dd1-markov}
Let $X \ge 0$ be a random variable with $\E[X] < \infty$. For every $t > 0$,
\begin{equation}
\label{eq:dd1-markov}
\Pr\bigl( X \ge t \bigr) \;\le\; \frac{\E[X]}{t}.
\end{equation}
\end{theorem}

\begin{proof}
Since $X \ge 0$, the pointwise inequality
$t\,\mathbf{1}\{X \ge t\} \le X$ holds on every outcome: if $X \ge t$ the left
side is $t \le X$; otherwise the left side is $0 \le X$. Taking expectations
and using monotonicity and linearity,
$t\,\Pr(X \ge t) = \E\bigl[t\,\mathbf{1}\{X \ge t\}\bigr] \le \E[X]$.
Divide by $t$.
\end{proof}

\begin{corollary}[Chebyshev's inequality]
\label{cor:dd1-chebyshev}
Let $X$ have mean $\mu$ and variance $\sigma^2 < \infty$. For every $t > 0$,
\begin{equation}
\label{eq:dd1-chebyshev}
\Pr\bigl( |X - \mu| \ge t \bigr) \;\le\; \frac{\sigma^2}{t^2}.
\end{equation}
\end{corollary}

\begin{proof}
Apply Theorem~\ref{thm:dd1-markov} to the nonnegative variable
$(X - \mu)^2$ at level $t^2$:
$\Pr(|X - \mu| \ge t) = \Pr\bigl((X - \mu)^2 \ge t^2\bigr)
\le \E[(X-\mu)^2] / t^2 = \sigma^2 / t^2$.
\end{proof}

Chebyshev is two-sided. When only one direction of deviation is harmful---a
niche head performing \emph{worse} on a fresh episode than on its stored
ones---we can do strictly better.

\begin{theorem}[Cantelli's one-sided inequality]
\label{thm:dd1-cantelli}
Let $X$ have mean $\mu$ and variance $\sigma^2 < \infty$. For every $t > 0$,
\begin{equation}
\label{eq:dd1-cantelli}
\Pr\bigl( X - \mu \ge t \bigr) \;\le\; \frac{\sigma^2}{\sigma^2 + t^2}.
\end{equation}
\end{theorem}

\begin{proof}
Let $Y = X - \mu$, so $\E[Y] = 0$ and $\E[Y^2] = \sigma^2$. Fix an arbitrary
auxiliary constant $u > 0$. On the event $\{Y \ge t\}$ we have
$Y + u \ge t + u > 0$, hence $(Y + u)^2 \ge (t + u)^2$ there, and
$(Y+u)^2 \ge 0$ everywhere. Therefore
\[
\Pr( Y \ge t )
\;=\; \Pr\bigl( Y + u \ge t + u \bigr)
\;\le\; \Pr\bigl( (Y + u)^2 \ge (t + u)^2 \bigr)
\;\le\; \frac{\E\bigl[(Y + u)^2\bigr]}{(t + u)^2}
\;=\; \frac{\sigma^2 + u^2}{(t + u)^2},
\]
where the middle step is Markov (Theorem~\ref{thm:dd1-markov}) applied to the
nonnegative variable $(Y + u)^2$, and the last step expands
$\E[(Y+u)^2] = \E[Y^2] + 2u\,\E[Y] + u^2 = \sigma^2 + u^2$.

The bound holds for every $u > 0$, so we minimize the right-hand side. Let
$\varphi(u) = (\sigma^2 + u^2)/(t + u)^2$. Then
\[
\varphi'(u)
= \frac{2u\,(t+u)^2 - 2(t+u)(\sigma^2 + u^2)}{(t+u)^4}
= \frac{2\bigl[u(t + u) - (\sigma^2 + u^2)\bigr]}{(t+u)^3}
= \frac{2\,(ut - \sigma^2)}{(t+u)^3},
\]
which vanishes at $u^* = \sigma^2 / t$, is negative for $u < u^*$ and positive
for $u > u^*$; so $u^*$ is the global minimizer on $(0, \infty)$.
Substituting,
\[
\varphi(u^*)
= \frac{\sigma^2 + \sigma^4/t^2}{\bigl(t + \sigma^2/t\bigr)^2}
= \frac{\sigma^2\,(t^2 + \sigma^2)/t^2}{(t^2 + \sigma^2)^2 / t^2}
= \frac{\sigma^2}{\sigma^2 + t^2}. \qedhere
\]
\end{proof}

\begin{remark}[Why one-sidedness matters for the episode-transfer lemma]
\label{rem:dd1-onesided}
In the transfer argument of Part II, the certified quantity is the average
regret of a niche head over its stored episodes, and the failure event is
one-directional: the fresh-episode regret exceeding that average by more than
a tolerance $t$. A fresh episode on which the head does \emph{better} than
its certificate is a windfall, not a failure. Cantelli prices exactly the
harmful tail: at $t = \sigma$, Chebyshev gives the vacuous bound $1$ split
across two tails (effectively $1/2$ per side by symmetry heuristics, but with
no valid one-sided guarantee better than $1$), while Cantelli gives
$\sigma^2/(2\sigma^2) = 1/2$ \emph{as a theorem}, and at $t = 2\sigma$ gives
$1/5$ versus Chebyshev's two-sided $1/4$. Moreover Cantelli requires no
boundedness and no symmetry---only a variance---which suits regret
distributions that are skewed by construction (they live in $[0, B]$ and pile
up near $0$ for a well-adapted head).
\end{remark}

\subsection{Hoeffding's Inequality}

Variance-based bounds decay polynomially in $t$. For \emph{bounded} variables
---and the regret is bounded in $[0, B]$ by
Proposition~\ref{prop:dd1-regret-range}---exponential tails are available.
The engine is a bound on the moment generating function.

\begin{lemma}[Hoeffding's lemma; proof sketch]
\label{lem:dd1-hoeffding-lemma}
Let $X \in [a, b]$ almost surely, with $\E[X] = 0$. Then for all
$\lambda \in \R$,
\begin{equation}
\label{eq:dd1-hoeffding-lemma}
\E\bigl[ e^{\lambda X} \bigr] \;\le\;
\exp\!\Bigl( \frac{\lambda^2 (b - a)^2}{8} \Bigr).
\end{equation}
\end{lemma}

\begin{proof}[Proof sketch]
Define the cumulant generating function $\psi(\lambda) = \log \E[e^{\lambda
X}]$, which is finite for all $\lambda$ (boundedness) and twice
differentiable. Direct computation gives $\psi(0) = 0$ and
$\psi'(0) = \E[X] = 0$. The second derivative is
\[
\psi''(\lambda) = \E_\lambda[X^2] - \bigl(\E_\lambda[X]\bigr)^2
= \Var_\lambda(X),
\]
the variance of $X$ under the exponentially tilted law
$d\Pr_\lambda \propto e^{\lambda X}\, d\Pr$ (this is the step we do not carry
out in full detail; it is a one-line differentiation under the expectation,
legitimate by dominated convergence on a bounded variable). The tilted law is
still supported in $[a, b]$, and \emph{any} random variable $V \in [a, b]$
satisfies
\[
\Var(V) \;\le\; \E\Bigl[\Bigl(V - \tfrac{a + b}{2}\Bigr)^2\Bigr]
\;\le\; \Bigl(\tfrac{b - a}{2}\Bigr)^2,
\]
where the first inequality holds because the variance minimizes the expected
squared deviation over all centering constants, and the second because
$|V - (a+b)/2| \le (b-a)/2$ pointwise. Hence
$\psi''(\lambda) \le (b - a)^2 / 4$ for all $\lambda$, and Taylor's theorem
with the exact remainder gives, for some $\xi$ between $0$ and $\lambda$,
\[
\psi(\lambda) = \psi(0) + \lambda\,\psi'(0)
+ \tfrac{\lambda^2}{2}\,\psi''(\xi)
\;\le\; \frac{\lambda^2 (b-a)^2}{8}.
\]
Exponentiate.
\end{proof}

\begin{theorem}[Hoeffding's inequality]
\label{thm:dd1-hoeffding}
Let $X_1, \dots, X_n$ be independent with $X_i \in [a_i, b_i]$ almost surely,
and let $\bar{X} = \frac{1}{n}\sum_{i=1}^n X_i$, $\mu = \E[\bar{X}]$. Then for
every $t > 0$,
\begin{equation}
\label{eq:dd1-hoeffding}
\Pr\bigl( \bar{X} - \mu \ge t \bigr)
\;\le\; \exp\!\Bigl( - \frac{2 n^2 t^2}{\sum_{i=1}^n (b_i - a_i)^2} \Bigr),
\end{equation}
and the same bound holds for $\Pr(\mu - \bar{X} \ge t)$; a union bound gives
the two-sided version with an extra factor $2$. In particular, if
$X_i \in [0, B]$ for all $i$, then with probability at least $1 - \delta$,
\begin{equation}
\label{eq:dd1-hoeffding-radius}
\bigl| \bar{X} - \mu \bigr| \;\le\; B \sqrt{ \frac{\ln(2/\delta)}{2n} }.
\end{equation}
\end{theorem}

\begin{proof}
Let $S = \sum_i (X_i - \E X_i)$, so $\{\bar{X} - \mu \ge t\} =
\{S \ge nt\}$. For any $\lambda > 0$, by Markov's inequality applied to the
nonnegative variable $e^{\lambda S}$ (the Chernoff step) and independence,
\[
\Pr( S \ge nt ) \;\le\; e^{-\lambda n t}\, \E\bigl[ e^{\lambda S} \bigr]
\;=\; e^{-\lambda n t} \prod_{i=1}^n \E\bigl[ e^{\lambda (X_i - \E X_i)}
\bigr]
\;\le\; \exp\!\Bigl( -\lambda n t + \frac{\lambda^2}{8} \sum_{i=1}^n
(b_i - a_i)^2 \Bigr),
\]
using Lemma~\ref{lem:dd1-hoeffding-lemma} for each centered $X_i - \E X_i \in
[a_i - \E X_i,\, b_i - \E X_i]$, an interval of the same width $b_i - a_i$.
The exponent $-\lambda n t + \frac{\lambda^2}{8}\sum_i (b_i - a_i)^2$ is a
convex parabola in $\lambda$, minimized at
$\lambda^* = 4nt / \sum_i (b_i - a_i)^2 > 0$, with minimum value
$-2 n^2 t^2 / \sum_i (b_i - a_i)^2$. This proves \eqref{eq:dd1-hoeffding};
the lower tail follows by applying the result to $-X_i$. For
\eqref{eq:dd1-hoeffding-radius}, set $a_i = 0$, $b_i = B$, equate the
two-sided bound $2\exp(-2nt^2/B^2)$ to $\delta$, and solve for $t$.
\end{proof}

\subsection{Empirical Bernstein: Paying the Observed Variance, Not the Range}

Hoeffding's radius \eqref{eq:dd1-hoeffding-radius} scales with the worst-case
range $B$, no matter how concentrated the data actually are. When the sample
is tiny this is painful, because $B$ for the decision regret is a gross
worst-case constant (Proposition~\ref{prop:dd1-regret-range}). Bernstein-type
inequalities replace $B$ by the standard deviation, but the classical
Bernstein bound needs the \emph{true} variance, which is unknown. The
empirical-Bernstein inequality of Maurer and Pontil
\cite{maurer2009empirical} closes the loop by using the \emph{sample}
variance, at the price of a lower-order $B/n$ term.

\begin{theorem}[Empirical Bernstein inequality \cite{maurer2009empirical};
stated without proof]
\label{thm:dd1-empbernstein}
Let $X_1, \dots, X_n$ ($n \ge 2$) be i.i.d.\ with values in $[0, B]$, sample
mean $\bar{X}$, and sample variance
$V_n = \frac{1}{n - 1}\sum_{i=1}^n (X_i - \bar{X})^2$. Then with probability
at least $1 - \delta$,
\begin{equation}
\label{eq:dd1-empbernstein}
\mu \;\le\; \bar{X} \;+\; \sqrt{ \frac{2\, V_n \ln(2/\delta)}{n} }
\;+\; \frac{7\, B \ln(2/\delta)}{3\,(n - 1)} .
\end{equation}
\end{theorem}

\begin{intuition}
Read \eqref{eq:dd1-empbernstein} as: \emph{you pay the spread you actually
observed} (the $\sqrt{V_n/n}$ term), plus a \emph{fixed entrance fee}
proportional to $B/n$ for the privilege of estimating the variance from the
same data. When the observations are tightly clustered ($\sqrt{V_n} \ll B$)
and $n$ is moderate, the first term is far below Hoeffding's
$B\sqrt{\ln(2/\delta)/2n}$ and the entrance fee is negligible. When $n$ is
very small, the entrance fee dominates everything.
\end{intuition}

\begin{remark}[Why this matters at $E_r = 5$--$7$ episodes]
\label{rem:dd1-small-n}
In \nm{}, the per-niche certificate averages the episode risks of the
$E_r \in \{5, \dots, 7\}$ episodes stored in niche $r$'s memory. The entire
design of the niche structure---episodes are routed to a niche precisely
because the niche head's decision behavior on them is similar---aims to make
the \emph{within-niche} spread of episode risks small, i.e.\ $V_n \ll B^2$.
Empirical Bernstein is the inequality that can, in principle, convert that
design property into a tighter certificate, because it is the only one in
this family whose radius adapts to the realized within-niche concentration.
Hoeffding cannot see the clustering at all. The sobering quantitative reality
at $n \in \{5, \dots, 7\}$ is worked out in the next example.
\end{remark}

\begin{example}[Estimation radii with $6$ episodes]
\label{ex:dd1-radius}
Take $n = 6$ episodes, regret range $B = 0.4$, confidence $\delta = 0.1$, so
$\ln(2/\delta) = \ln 20 \approx 2.9957$.

\emph{Hoeffding.} The two-sided radius \eqref{eq:dd1-hoeffding-radius} is
\[
B \sqrt{\frac{\ln(2/\delta)}{2n}}
= 0.4 \times \sqrt{\frac{2.9957}{12}}
= 0.4 \times \sqrt{0.24964}
= 0.4 \times 0.49964
\approx 0.1999 \approx 0.20 .
\]
The certificate is ``mean regret within $\pm 0.20$''---\emph{half the entire
range} $[0, 0.4]$. As a quantitative guarantee this is nearly vacuous.

\emph{Empirical Bernstein.} Suppose the six stored episode regrets are
tightly clustered with sample standard deviation $\sqrt{V_n} = 0.05$. The
variance term of \eqref{eq:dd1-empbernstein} is
\[
\sqrt{\frac{2 \times 0.0025 \times 2.9957}{6}}
= \sqrt{0.0024964} \approx 0.050,
\]
a four-fold improvement---but the entrance fee is
\[
\frac{7 \times 0.4 \times 2.9957}{3 \times 5}
= \frac{8.388}{15} \approx 0.559,
\]
which alone exceeds the whole range $B$. At $n = 6$ the empirical-Bernstein
certificate is \emph{worse} than Hoeffding's; the variance adaptivity only
pays once $n$ is large enough that the $B/(n-1)$ term recedes (here, roughly
$n \gtrsim 40$ for these constants).
\end{example}

\begin{warning}
Example~\ref{ex:dd1-radius} is the honest small-sample accounting that frames
every finite-sample statement in this document: at $E_r = 5$--$7$ episodes,
\emph{no} distribution-free concentration inequality certifies the per-niche
mean regret to a useful tolerance. The role of these bounds in the \nm{}
theory is therefore \emph{structural}, not numerical: they identify
\emph{which} quantities control transfer (the within-niche mean and
variance of episode risk, and only those), and they predict the correct
\emph{qualitative} ordering of mechanisms (lower within-niche variance
$\Rightarrow$ tighter transfer). That prediction is exactly what the
variance-penalized objective of Section~\ref{sec:dd1-invariance} acts on, and
it is tested empirically rather than certified theoretically.
\end{warning}

% ============================================================================
\section{Distances Between Distributions and the Pinsker Bridge}
\label{sec:dd1-distances}

\nm{} inherits from the invariant predict-and-optimize line of work
\cite{yan2026invpnco} guarantees of the form ``within a regime, the
divergence between the training mixture of episodes and a deployment episode
is at most $\epsilon$.'' Such statements are about \emph{distributions}; what
we need are statements about \emph{regret}. This section builds the standard
bridge: total variation converts distributional closeness into closeness of
bounded expectations, and Pinsker's inequality converts KL divergence into
total variation.

Throughout, $P$ and $Q$ are probability measures on a common measurable space
$(\Omega, \mathcal{F})$, and we fix a dominating $\sigma$-finite measure
$\mu$ (e.g.\ $\mu = P + Q$) with densities $p = dP/d\mu$, $q = dQ/d\mu$. All
statements are independent of the choice of $\mu$.

\subsection{Total Variation Distance}

\begin{definition}[Total variation]
\label{def:dd1-tv}
The \textbf{total variation distance} between $P$ and $Q$ is
\begin{equation}
\label{eq:dd1-tv-def}
\TV(P, Q) \;=\; \sup_{A \in \mathcal{F}} \bigl| P(A) - Q(A) \bigr|.
\end{equation}
\end{definition}

\begin{proposition}[Variational and $L_1$ characterizations]
\label{prop:dd1-tv-variational}
With $p, q$ as above,
\begin{equation}
\label{eq:dd1-tv-three}
\TV(P, Q)
\;=\; \frac{1}{2} \int_\Omega |p - q| \, d\mu
\;=\; \sup_{g : \Omega \to [0, 1]} \Bigl| \E_P[g] - \E_Q[g] \Bigr|,
\end{equation}
where the supremum is over all measurable functions $g$ with values in
$[0,1]$, and is attained at $g = \mathbf{1}_{A^*}$ with
$A^* = \{p \ge q\}$.
\end{proposition}

\begin{proof}
Write $\Delta = p - q$ and note the key cancellation
\begin{equation}
\label{eq:dd1-cancellation}
\int_\Omega \Delta \, d\mu = \int p \, d\mu - \int q \, d\mu = 1 - 1 = 0
\quad\Longrightarrow\quad
\int_{A^*} \Delta \, d\mu \;=\; \int_{(A^*)^c} (-\Delta) \, d\mu ,
\end{equation}
where $A^* = \{\Delta \ge 0\}$; both sides of the implication equal
$\frac{1}{2}\int |\Delta| \, d\mu$ since
$\int |\Delta| = \int_{A^*} \Delta + \int_{(A^*)^c} (-\Delta)$.

\emph{Step 1: every $[0,1]$-valued $g$ is dominated.} For any such $g$,
\[
\E_P[g] - \E_Q[g] = \int g\, \Delta \, d\mu
= \int_{A^*} g\, \Delta \, d\mu + \int_{(A^*)^c} g\, \Delta \, d\mu
\;\le\; \int_{A^*} 1 \cdot \Delta \, d\mu + \int_{(A^*)^c} 0 \cdot \Delta \,
d\mu
= \tfrac{1}{2} \textstyle\int |\Delta| \, d\mu,
\]
using $0 \le g \le 1$, $\Delta \ge 0$ on $A^*$, and $\Delta < 0$ on
$(A^*)^c$. Applying the same bound to $1 - g$ (also $[0,1]$-valued) gives
$\E_Q[g] - \E_P[g] \le \frac{1}{2}\int|\Delta|$ as well (note
$\E_Q[g]-\E_P[g] = \E_P[1-g]-\E_Q[1-g]$). Hence
$\sup_g |\E_P g - \E_Q g| \le \frac{1}{2}\int |\Delta|\, d\mu$.

\emph{Step 2: attainment.} The choice $g = \mathbf{1}_{A^*}$ achieves
$\E_P[g] - \E_Q[g] = \int_{A^*} \Delta\, d\mu = \frac{1}{2}\int|\Delta|\,
d\mu$ by \eqref{eq:dd1-cancellation}. So the supremum over $[0,1]$-valued $g$
\emph{equals} $\frac{1}{2}\int |\Delta|$.

\emph{Step 3: events.} Indicators $\mathbf{1}_A$ are particular
$[0,1]$-valued functions, so
$\sup_A |P(A) - Q(A)| \le \sup_g |\E_P g - \E_Q g|$; and the optimizer found
in Step 2 \emph{is} an indicator, so the two suprema coincide, and both equal
$\frac{1}{2}\int|p - q|\,d\mu$.
\end{proof}

\begin{intuition}
$\TV$ answers: \emph{over all bounded tests normalized to $[0,1]$, how
differently can the two distributions score?} The optimal test is the
likelihood-ratio region $\{p \ge q\}$. The $[0,1]$-function form---rather
than the events-only form---is the one we will actually use, because a
bounded loss rescaled into $[0,1]$ is such a function but is not an
indicator.
\end{intuition}

\subsection{KL Divergence and Gibbs' Inequality}

\begin{definition}[Kullback--Leibler divergence]
\label{def:dd1-kl}
If $P \ll Q$ (i.e.\ $q = 0 \Rightarrow p = 0$, $\mu$-a.e.),
\begin{equation}
\label{eq:dd1-kl-def}
\KL(P \,\|\, Q) \;=\; \int_{\{p > 0\}} p \,\log \frac{p}{q} \, d\mu
\;=\; \E_P\Bigl[ \log \frac{p}{q} \Bigr],
\end{equation}
with the conventions $0 \log 0 = 0$; if $P \not\ll Q$ we set
$\KL(P\|Q) = +\infty$.
\end{definition}

\begin{proposition}[Gibbs' inequality / non-negativity]
\label{prop:dd1-gibbs}
$\KL(P \| Q) \ge 0$, with equality if and only if $P = Q$.
\end{proposition}

\begin{proof}
Assume $P \ll Q$ (otherwise $\KL = \infty > 0$). Then
\[
-\KL(P\|Q) = \E_P\Bigl[ \log \frac{q}{p} \Bigr]
\;\le\; \log \E_P\Bigl[ \frac{q}{p} \Bigr]
= \log \int_{\{p > 0\}} q \, d\mu
\;\le\; \log 1 = 0,
\]
where the first inequality is Jensen's inequality for the concave function
$\log$, and the second uses $\int_{\{p>0\}} q\, d\mu \le \int q \, d\mu = 1$.
For the equality case: $\log$ is \emph{strictly} concave, so Jensen is an
equality iff $q/p$ is $P$-a.s.\ constant; combined with
$\int_{\{p>0\}} q = 1$ (forced by the second inequality being tight) the
constant is $1$, i.e.\ $p = q$ $\mu$-a.e., i.e.\ $P = Q$.
\end{proof}

\subsection{Pinsker's Inequality}

We give the classical proof \cite{tsybakov2009introduction} in two moves: a
data-processing (two-point) reduction via the log-sum inequality, then the
scalar inequality for Bernoulli pairs.

\begin{lemma}[Log-sum inequality]
\label{lem:dd1-logsum}
For nonnegative reals $a_1, \dots, a_m$ and positive reals
$b_1, \dots, b_m$,
\begin{equation}
\label{eq:dd1-logsum}
\sum_{i=1}^m a_i \log \frac{a_i}{b_i}
\;\ge\; \Bigl( \sum_{i=1}^m a_i \Bigr)
\log \frac{\sum_{i=1}^m a_i}{\sum_{i=1}^m b_i},
\end{equation}
with the convention $0 \log 0 = 0$. The same holds with sums replaced by
integrals of nonnegative functions.
\end{lemma}

\begin{proof}
Let $\phi(t) = t \log t$ for $t > 0$, $\phi(0) = 0$; $\phi$ is convex on
$[0, \infty)$ since $\phi''(t) = 1/t > 0$. Set $b = \sum_i b_i$ and apply
Jensen's inequality with weights $b_i / b$ to the points $t_i = a_i / b_i$:
\[
\sum_{i=1}^m \frac{b_i}{b}\, \phi\Bigl( \frac{a_i}{b_i} \Bigr)
\;\ge\; \phi\Bigl( \sum_{i=1}^m \frac{b_i}{b} \cdot \frac{a_i}{b_i} \Bigr)
= \phi\Bigl( \frac{\sum_i a_i}{b} \Bigr).
\]
Multiplying both sides by $b$ and expanding $\phi$:
\[
\sum_i b_i \cdot \frac{a_i}{b_i} \log\frac{a_i}{b_i}
= \sum_i a_i \log \frac{a_i}{b_i}
\;\ge\; b \cdot \frac{\sum_i a_i}{b} \log \frac{\sum_i a_i}{b}
= \Bigl(\sum_i a_i\Bigr) \log \frac{\sum_i a_i}{\sum_i b_i}.
\]
The integral version is Jensen for the same $\phi$ with the normalized
measure $b(\omega)\,d\mu / \int b \, d\mu$.
\end{proof}

\begin{lemma}[Two-point (data-processing) reduction]
\label{lem:dd1-dataproc}
For any event $A \in \mathcal{F}$, writing $a = P(A)$, $b = Q(A)$ and
$\mathrm{kl}(a \,\|\, b) = a \log\frac{a}{b} + (1-a)\log\frac{1-a}{1-b}$ for
the Bernoulli KL,
\begin{equation}
\label{eq:dd1-dataproc}
\KL(P \,\|\, Q) \;\ge\; \mathrm{kl}\bigl( P(A) \,\|\, Q(A) \bigr).
\end{equation}
\end{lemma}

\begin{proof}
Split the defining integral over the partition $\{A, A^c\}$ and apply the
integral log-sum inequality (Lemma~\ref{lem:dd1-logsum}) on each cell:
\[
\KL(P\|Q)
= \int_A p \log \frac{p}{q}\, d\mu + \int_{A^c} p \log \frac{p}{q}\, d\mu
\;\ge\; \Bigl(\int_A p\Bigr) \log \frac{\int_A p}{\int_A q}
+ \Bigl(\int_{A^c} p\Bigr) \log \frac{\int_{A^c} p}{\int_{A^c} q}
= \mathrm{kl}\bigl(P(A) \| Q(A)\bigr). \qedhere
\]
\end{proof}

\begin{lemma}[Scalar Pinsker]
\label{lem:dd1-scalar-pinsker}
For all $a, b \in [0, 1]$ (with $\mathrm{kl}(a\|b) = \infty$ when undefined),
\begin{equation}
\label{eq:dd1-scalar-pinsker}
\mathrm{kl}(a \,\|\, b) \;\ge\; 2\,(a - b)^2 .
\end{equation}
\end{lemma}

\begin{proof}
If $b \in \{0, 1\}$ and $a \neq b$ the left side is $+\infty$ and the claim
is trivial; if $a = b$ both sides vanish. So fix $b \in (0,1)$ and define on
$(0,1)$ (extending continuously to the endpoints where finite)
\[
f(a) \;=\; \mathrm{kl}(a \| b) - 2 (a - b)^2
\;=\; a \log\frac{a}{b} + (1 - a)\log\frac{1-a}{1-b} - 2(a-b)^2 .
\]
Then $f(b) = 0$, and differentiating in $a$:
\[
f'(a) = \log \frac{a}{b} - \log \frac{1-a}{1-b} - 4(a - b),
\qquad f'(b) = 0,
\]
\[
f''(a) = \frac{1}{a} + \frac{1}{1 - a} - 4 = \frac{1}{a(1-a)} - 4 \;\ge\; 0,
\]
since $a(1-a) \le 1/4$ for $a \in (0,1)$ (AM--GM: $a(1-a) \le
\bigl(\frac{a + (1-a)}{2}\bigr)^2 = \frac14$). Thus $f'$ is nondecreasing
and vanishes at $a = b$: $f' \le 0$ on $(0, b]$ and $f' \ge 0$ on $[b, 1)$,
so $f$ attains its minimum over $(0,1)$ at $a = b$, where it is $0$. Hence
$f(a) \ge 0$ for all $a$, which is \eqref{eq:dd1-scalar-pinsker}.
\end{proof}

\begin{theorem}[Pinsker's inequality]
\label{thm:dd1-pinsker}
For all probability measures $P, Q$,
\begin{equation}
\label{eq:dd1-pinsker}
\TV(P, Q) \;\le\; \sqrt{ \tfrac{1}{2}\, \KL(P \,\|\, Q) } .
\end{equation}
\end{theorem}

\begin{proof}
If $\KL(P\|Q) = \infty$ there is nothing to prove. Otherwise let
$A^* = \{p \ge q\}$, so that by
Proposition~\ref{prop:dd1-tv-variational}, $\TV(P, Q) = P(A^*) - Q(A^*)$.
Chaining Lemmas~\ref{lem:dd1-dataproc} and \ref{lem:dd1-scalar-pinsker},
\[
\KL(P \| Q)
\;\ge\; \mathrm{kl}\bigl(P(A^*) \,\|\, Q(A^*)\bigr)
\;\ge\; 2\,\bigl(P(A^*) - Q(A^*)\bigr)^2
\;=\; 2\, \TV(P, Q)^2 .
\]
Rearranging gives \eqref{eq:dd1-pinsker}.
\end{proof}

\subsection{The Pinsker Bridge: from Divergence to Regret}

\begin{theorem}[Bounded-loss transfer / Pinsker bridge]
\label{thm:dd1-bridge}
Let $\ell : \Omega \to [0, B]$ be any measurable loss, and write
$L_P = \E_P[\ell]$, $L_Q = \E_Q[\ell]$. Then
\begin{equation}
\label{eq:dd1-bridge}
\bigl| L_P - L_Q \bigr|
\;\le\; B \cdot \TV(P, Q)
\;\le\; B \, \sqrt{ \tfrac{1}{2} \,\KL(P \,\|\, Q) } .
\end{equation}
\end{theorem}

\begin{proof}
The function $g = \ell / B$ takes values in $[0, 1]$, so by the variational
characterization (Proposition~\ref{prop:dd1-tv-variational}),
\[
\bigl| L_P - L_Q \bigr| = B \,\bigl| \E_P[g] - \E_Q[g] \bigr|
\;\le\; B \sup_{g' : 0 \le g' \le 1} \bigl| \E_P[g'] - \E_Q[g'] \bigr|
= B \cdot \TV(P, Q).
\]
The second inequality in \eqref{eq:dd1-bridge} is Pinsker
(Theorem~\ref{thm:dd1-pinsker}).
\end{proof}

\begin{keypoint}
\textbf{Role in the \nm{} theory.} Fix a regime $r$ and apply
Theorem~\ref{thm:dd1-bridge} with: $\Omega$ the space of (feature, return)
days; $\ell(\omega) = \rho(\hat{y}_\theta(x_\omega); y_\omega)$ the decision
regret of niche head $\theta$, bounded by $B = 2 k\, w_{\max}\, y_{\max}$
(Proposition~\ref{prop:dd1-regret-range}); $P$ the mixture of the niche's
stored training episodes; $Q$ a fresh deployment episode of the same regime.
An Inv-PnCO-style assumption \cite{yan2026invpnco}---that episodes of one
regime are divergence-close, $\KL(Q \| P) \le \epsilon_r$ (or
$\TV \le \tau_r$)---then yields a \emph{regret} transfer guarantee:
\[
\text{fresh-episode regret}
\;\le\; \text{stored-episode regret} + B\sqrt{\epsilon_r / 2},
\]
which composes with the concentration bounds of
Section~\ref{sec:dd1-concentration} (controlling the stored-episode average)
to give the episode-transfer lemma of Part II. The bridge is what lets a
statement about \emph{data distributions} (checkable by regime statistics)
control a statement about \emph{money} (regret of induced portfolios), for
\emph{any} head simultaneously---no union bound over heads is needed because
\eqref{eq:dd1-bridge} is deterministic given the divergence.
\end{keypoint}

\begin{remark}[Honest looseness accounting]
\label{rem:dd1-B-loose}
Every appearance of $B = 2 k\, w_{\max}\, y_{\max}$ in
\eqref{eq:dd1-bridge} imports the worst case of
Proposition~\ref{prop:dd1-regret-range}: the forecast shorting itself into
the $k$ worst assets while the oracle rides the $k$ best, all at the extreme
return $y_{\max}$. Empirically observed decision regrets in the companion
experiments range roughly \emph{fifty times smaller} than $B$. The bridge's
\emph{form}---transfer error scaling as
$(\text{regret scale}) \times \sqrt{\text{divergence}}$---is the content we
rely on; its \emph{constant} should be read as a certificate of structure,
not of magnitude. Tightening $B$ to a high-probability regret envelope (e.g.\
a quantile of the realized regret distribution) is possible but converts the
deterministic bridge into a probabilistic one; we keep the loose constant and
flag it.
\end{remark}

% ============================================================================
\section{Online Learning: Hedge, Exponentiated Gradient, and Their Limits}
\label{sec:dd1-online}

The \nm{} gate maintains a probability vector over niches and updates it
multiplicatively in observed decision quality. This update has three classical
pedigrees---Hedge \cite{freund1997decision}, exponentiated gradient
\cite{kivinen1997exponentiated}, and multiplicative weights
\cite{arora2012multiplicative}---and one classical guarantee, the
$O(\sqrt{T \log N})$ adversarial regret bound \cite{cesa2006prediction}. This
section develops all of it, and then states precisely why the classical
guarantee does \emph{not} cover \nm{}'s use of the update.

\subsection{The Experts Problem and the Hedge Update}

\begin{definition}[Prediction with expert advice]
\label{def:dd1-experts}
There are $N$ \textbf{experts}. At each round $t = 1, \dots, T$: the learner
chooses a distribution $p_t \in \Delta_N = \{p \in [0,1]^N : \sum_i p_i =
1\}$; an adversary simultaneously chooses a loss vector
$\ell_t \in [0, 1]^N$; the learner suffers $\langle p_t, \ell_t \rangle$ and
observes all of $\ell_t$. The learner's \textbf{regret} is
\begin{equation}
\label{eq:dd1-experts-regret}
R_T \;=\; \sum_{t=1}^T \langle p_t, \ell_t \rangle
\;-\; \min_{i \in [N]} \sum_{t=1}^T \ell_{t,i},
\end{equation}
the gap to the best \emph{single} expert in hindsight.
\end{definition}

\begin{definition}[Hedge / multiplicative weights]
\label{def:dd1-hedge}
Fix a learning rate $\eta > 0$. Initialize $w_{1,i} = 1$ for all $i$ and
$p_1 = (1/N, \dots, 1/N)$. After observing $\ell_t$,
\begin{equation}
\label{eq:dd1-hedge}
w_{t+1, i} = w_{t, i}\, e^{-\eta\, \ell_{t,i}},
\qquad
p_{t+1, i} = \frac{w_{t+1, i}}{\sum_{j=1}^N w_{t+1, j}}
= \frac{p_{t,i}\, e^{-\eta \ell_{t,i}}}{\sum_j p_{t,j}\, e^{-\eta
\ell_{t,j}}}.
\end{equation}
\end{definition}

This lineage runs through universal portfolio selection
\cite{cover1991universal} and online portfolio optimization
\cite{li2014online}: there, the ``experts'' are constant-rebalanced
portfolios or trading strategies and the weights track wealth shares. In
\nm{} the experts are niche heads and the weights drive both selection and
training-resource allocation---a difference we return to in
Section~\ref{sec:dd1-online-limits}.

\subsection{Exponentiated Gradient as Mirror Descent with Entropy}

The exponential form \eqref{eq:dd1-hedge} is not an ad hoc choice: it is the
\emph{exact} solution of a regularized linear step where the regularizer is
the KL divergence on the simplex \cite{kivinen1997exponentiated}.

\begin{proposition}[EG update solves the KL-regularized linear step]
\label{prop:dd1-eg-mirror}
Fix $p_t$ in the relative interior of $\Delta_N$ and $\eta > 0$. The unique
solution of
\begin{equation}
\label{eq:dd1-mirror-step}
p_{t+1} \;=\; \argmin_{p \in \Delta_N}
\Bigl\{ \eta\, \langle \ell_t, p \rangle + \KL(p \,\|\, p_t) \Bigr\},
\qquad
\KL(p \| p_t) = \sum_{i=1}^N p_i \log \frac{p_i}{p_{t,i}},
\end{equation}
is exactly the Hedge update \eqref{eq:dd1-hedge}.
\end{proposition}

\begin{proof}
The objective $J(p) = \eta \langle \ell_t, p\rangle + \KL(p\|p_t)$ is
strictly convex on $\Delta_N$ ($\langle \ell_t, \cdot\rangle$ is linear and
$\KL(\cdot \| p_t)$ is strictly convex, being a sum of strictly convex scalar
functions $p_i \mapsto p_i \log(p_i / p_{t,i})$), so the minimizer is unique.
Moreover the entropy term's derivative blows up at the boundary
($\partial_{p_i} p_i \log p_i = \log p_i + 1 \to -\infty$ as $p_i \to 0^+$),
which forces the minimizer into the relative interior; hence the
nonnegativity constraints are inactive and we may use a Lagrangian with only
the equality constraint. Set
\[
\mathcal{L}(p, \lambda)
= \sum_{i=1}^N \Bigl[ \eta\, \ell_{t,i}\, p_i
+ p_i \log \frac{p_i}{p_{t,i}} \Bigr]
+ \lambda \Bigl( \sum_{i=1}^N p_i - 1 \Bigr).
\]
Stationarity in each coordinate:
\[
\frac{\partial \mathcal{L}}{\partial p_i}
= \eta\, \ell_{t,i} + \log \frac{p_i}{p_{t,i}} + 1 + \lambda = 0
\quad\Longleftrightarrow\quad
p_i = p_{t,i}\, e^{-\eta \ell_{t,i}}\, e^{-(1 + \lambda)} .
\]
The factor $e^{-(1+\lambda)}$ is a positive constant independent of $i$;
the constraint $\sum_i p_i = 1$ pins it to
$\bigl(\sum_j p_{t,j} e^{-\eta \ell_{t,j}}\bigr)^{-1}$. This is precisely
\eqref{eq:dd1-hedge}.
\end{proof}

\begin{intuition}
The step \eqref{eq:dd1-mirror-step} balances two pulls: ``move mass toward
low-loss coordinates'' (the linear term) against ``do not move far from where
you are, measured \emph{multiplicatively}'' (the KL term). Measuring movement
by KL rather than Euclidean distance is what makes the update respect the
simplex geometry: probabilities can shrink toward zero but never jump across
it, weights never go negative, and a coordinate's \emph{relative} change
$p_{t+1,i}/p_{t,i}$ depends only on its own loss and the normalizer. This is
the precise sense in which the \nm{} gate is a ``well-behaved simplex
ascent'': bounded relative updates, invariance to common loss shifts, and no
projection step that could zero out a niche in one round.
\end{intuition}

\subsection{The Classical Regret Guarantee}

\begin{theorem}[Hedge regret bound \cite{freund1997decision,
cesa2006prediction}]
\label{thm:dd1-hedge-regret}
For losses $\ell_t \in [0,1]^N$ and any $\eta > 0$, Hedge guarantees
\begin{equation}
\label{eq:dd1-hedge-bound}
R_T \;\le\; \frac{\ln N}{\eta} + \frac{\eta\, T}{8},
\end{equation}
and with the tuned rate $\eta = \sqrt{8 \ln N / T}$,
$R_T \le \sqrt{(T/2) \ln N} = O(\sqrt{T \log N})$.
\end{theorem}

\begin{proof}[Proof (potential-function argument)]
Let $W_t = \sum_{i=1}^N w_{t,i}$ with the weights of
Definition~\ref{def:dd1-hedge}, and consider the potential $\ln W_t$.

\emph{Per-round drop.} Since $p_{t,i} = w_{t,i} / W_t$,
\[
\ln \frac{W_{t+1}}{W_t}
= \ln \sum_{i=1}^N \frac{w_{t,i}}{W_t}\, e^{-\eta \ell_{t,i}}
= \ln \E_{i \sim p_t}\bigl[ e^{-\eta \ell_{t,i}} \bigr]
\;\le\; -\eta\, \E_{i \sim p_t}[\ell_{t,i}] + \frac{\eta^2}{8}
= -\eta\, \langle p_t, \ell_t \rangle + \frac{\eta^2}{8},
\]
where the inequality is Hoeffding's lemma
(Lemma~\ref{lem:dd1-hoeffding-lemma}) applied to the random variable
$-\eta\,(\ell_{t,I} - \E\ell_{t,I})$, $I \sim p_t$, supported on an interval
of width $\eta \cdot 1$. Summing over $t$ and telescoping,
\begin{equation}
\label{eq:dd1-potential-upper}
\ln W_{T+1} - \ln W_1
\;\le\; -\eta \sum_{t=1}^T \langle p_t, \ell_t \rangle + \frac{\eta^2 T}{8}.
\end{equation}

\emph{Comparator floor.} For any fixed expert $i^*$, since all weights are
positive,
\begin{equation}
\label{eq:dd1-potential-lower}
\ln W_{T+1} \;\ge\; \ln w_{T+1, i^*}
= \ln\Bigl( 1 \cdot \prod_{t=1}^T e^{-\eta \ell_{t, i^*}} \Bigr)
= -\eta \sum_{t=1}^T \ell_{t, i^*},
\end{equation}
while $\ln W_1 = \ln N$.

\emph{Combine.} Chaining \eqref{eq:dd1-potential-lower} $\le \ln W_{T+1}$
with \eqref{eq:dd1-potential-upper}:
\[
-\eta \sum_t \ell_{t, i^*} - \ln N
\;\le\; -\eta \sum_t \langle p_t, \ell_t \rangle + \frac{\eta^2 T}{8}
\;\;\Longrightarrow\;\;
\sum_t \langle p_t, \ell_t \rangle - \sum_t \ell_{t, i^*}
\;\le\; \frac{\ln N}{\eta} + \frac{\eta T}{8}.
\]
Taking $i^*$ to be the best expert gives \eqref{eq:dd1-hedge-bound};
substituting the tuned $\eta$ gives
$2\sqrt{(\ln N)(T)/8} = \sqrt{(T/2)\ln N}$.
\end{proof}

Remarkably, the bound holds against an adversary who sees the learner's
distribution before choosing losses---but, crucially, an adversary whose
losses are \emph{handed to fixed experts}, a point we scrutinize below.

\subsection{The Replicator Limit}

For small $\eta$, the discrete multiplicative update becomes the replicator
dynamics of evolutionary game theory---the lens through which the
winner-take-all niche competition of \nm{} (and of GAUSE before it
\cite{li2026gause}) is analyzed.

\begin{proposition}[Small-$\eta$ replicator limit]
\label{prop:dd1-replicator}
For the Hedge update \eqref{eq:dd1-hedge} with bounded losses,
\begin{equation}
\label{eq:dd1-replicator}
p_{t+1, i} - p_{t, i}
\;=\; \eta\, p_{t,i} \bigl( \langle p_t, \ell_t \rangle - \ell_{t,i} \bigr)
\;+\; O(\eta^2),
\end{equation}
so as $\eta \to 0$ the rescaled-time trajectories solve the replicator ODE
$\dot{p}_i = p_i (\bar{f} - f_i)$ with ``fitness'' $f_i = \ell_{t,i}$ (lower
loss = higher fitness, $\bar{f} = \langle p, \ell\rangle$ the population
mean).
\end{proposition}

\begin{proof}
Expand the exponentials to first order: $e^{-\eta \ell_{t,i}} = 1 - \eta
\ell_{t,i} + O(\eta^2)$, uniformly in $i$ by boundedness. The normalizer is
\[
\sum_j p_{t,j} e^{-\eta \ell_{t,j}}
= 1 - \eta \langle p_t, \ell_t \rangle + O(\eta^2),
\qquad\text{so}\qquad
\Bigl( \sum_j p_{t,j} e^{-\eta \ell_{t,j}} \Bigr)^{-1}
= 1 + \eta \langle p_t, \ell_t \rangle + O(\eta^2).
\]
Multiplying,
\[
p_{t+1,i} = p_{t,i} \bigl(1 - \eta \ell_{t,i}\bigr)
\bigl(1 + \eta \langle p_t, \ell_t\rangle\bigr) + O(\eta^2)
= p_{t,i}\Bigl[ 1 + \eta\bigl( \langle p_t, \ell_t \rangle - \ell_{t,i}
\bigr) \Bigr] + O(\eta^2). \qedhere
\]
\end{proof}

\begin{intuition}
A niche grows in gate weight iff it beats the \emph{current
weighted average} of all niches; the growth rate is proportional to both its
advantage and its current weight. This is exactly the dynamics under which
competitive exclusion and niche differentiation are usually derived---and it
is a \emph{dynamical-systems} object (fixed points, basins, exclusion), not a
regret object.
\end{intuition}

\subsection{What the Classical Guarantee Does Not Cover}
\label{sec:dd1-online-limits}

\begin{warning}
\textbf{The adversarial regret bound does not transfer to \nm{}.}
Theorem~\ref{thm:dd1-hedge-regret} compares the learner to the best
\emph{fixed} expert under loss sequences that are \emph{exogenous}: expert
$i$'s loss at round $t$ does not depend on how much weight the learner has
been giving expert $i$. In \nm{} both premises fail, for the same reason they
fail in GAUSE \cite{li2026gause}:
\begin{enumerate}
\item \emph{Losses are endogenous.} Niche $i$'s decision quality at round $t$
depends on its head parameters, which were trained with data and gradient
budget allocated \emph{in proportion to past gate weights}. The loss sequence
$\ell_{t,i}$ is a functional of the learner's own trajectory
$p_{1:t-1}$---outside the oblivious- and even the standard adaptive-adversary
model, where ``expert $i$'s loss'' must remain well-defined independently of
the comparison.
\item \emph{The comparator is meaningless.} ``The best fixed niche in
hindsight'' is not a coherent benchmark when the niche's identity (its
trained parameters) is itself an artifact of the weight history; under a
different weight history the ``same'' niche would be a different predictor.
A regret bound against an incoherent comparator is vacuous even if formally
true.
\item \emph{Winner-take-all coupling.} The \nm{} assignment concentrates
training on high-weight niches, which is the very mechanism that creates
specialization; this positive feedback is what the replicator analysis
(Proposition~\ref{prop:dd1-replicator}) studies and what a worst-case regret
analysis is designed to neutralize.
\end{enumerate}
Accordingly, \nm{} uses the EG/Hedge update for its \emph{structural}
virtues established in Proposition~\ref{prop:dd1-eg-mirror} and the
intuition following it---a strictly interior, multiplicatively bounded,
geometry-respecting ascent on the simplex---and \emph{not} for the bound
\eqref{eq:dd1-hedge-bound}. Every stability claim about the coupled dynamics
is proved separately in Part II; none is inherited from
Theorem~\ref{thm:dd1-hedge-regret}.
\end{warning}

\subsection{Sleeping Experts and Adaptive Regret: the Gap \nm{} Addresses}
\label{sec:dd1-sleeping}

Non-stationarity has its own classical formalisms, and it is important to
locate precisely where they stop short of the \nm{} setting.

\begin{definition}[Sleeping experts \cite{freund1997sleeping}]
\label{def:dd1-sleeping}
In the sleeping-experts (specialists) problem, at each round $t$ an
adversary additionally reveals an \textbf{awake set} $A_t \subseteq [N]$;
only awake experts make predictions, the learner must put its mass on $A_t$,
and asleep experts are unevaluated. The standard guarantee is per-expert: for
every expert $i$,
\begin{equation}
\label{eq:dd1-sleeping-regret}
\sum_{t \,:\, i \in A_t} \langle p_t, \ell_t \rangle
\;-\; \sum_{t \,:\, i \in A_t} \ell_{t,i}
\;\le\; O\bigl( \sqrt{ T_i \,\log N } \bigr),
\qquad T_i = |\{ t : i \in A_t \}|,
\end{equation}
i.e.\ the learner competes with each specialist \emph{on the rounds where
that specialist is awake}.
\end{definition}

\begin{definition}[Adaptive and strongly adaptive regret
\cite{hazan2009adaptive, daniely2015strongly}]
\label{def:dd1-adaptive}
The \textbf{adaptive regret} of an algorithm is its worst regret over any
contiguous interval,
\begin{equation}
\label{eq:dd1-adaptive-regret}
\mathrm{AR}_T \;=\; \max_{1 \le s \le e \le T}
\Bigl[ \sum_{t=s}^{e} \langle p_t, \ell_t \rangle
- \min_{i} \sum_{t=s}^{e} \ell_{t,i} \Bigr];
\end{equation}
an algorithm is \textbf{strongly adaptive} \cite{daniely2015strongly} if for
every interval length $\tau$ its regret on every length-$\tau$ interval is
$O(\mathrm{poly}(\log T) \cdot \sqrt{\tau})$. Such algorithms track the best
expert \emph{per regime segment} rather than globally.
\end{definition}

\begin{keypoint}
\textbf{Why dormant \emph{fixed} experts retain trivially.} In both
frameworks an expert is a fixed function (or a fixed external advisor). When
expert $i$ sleeps for a thousand rounds, nothing happens to it: its
parameters do not exist or do not change, and when it wakes its predictions
are exactly what they would have been. Retention of dormant competence is
\emph{free}---it is an assumption, not an achievement. The guarantees
\eqref{eq:dd1-sleeping-regret} and \eqref{eq:dd1-adaptive-regret} are
therefore statements purely about the \emph{weighting} layer: how fast the
learner re-concentrates on the right expert after a regime switch.
\end{keypoint}

\begin{remark}[The precise gap]
\label{rem:dd1-gap}
\nm{} operates in a setting these formalisms do not model:
\textbf{capacity-bounded \emph{learning} experts under reward-driven
training allocation}. Formally: each expert $i$ carries parameters
$\theta_{t,i}$ in a bounded-capacity class; at each round a shared training
budget is divided according to (a function of) the gate weights $p_t$, and
$\theta_{t+1, i}$ is produced by gradient steps on data weighted by expert
$i$'s allocation; the loss $\ell_{t,i} = \ell(\theta_{t,i};\,
\text{episode}_t)$ is thereby a functional of the entire weight history. In
this setting dormancy is no longer safe by assumption: an expert that
receives weight (hence budget, hence gradient updates on
\emph{out-of-niche} data, or no updates while the shared feature extractor
drifts) can have its dormant competence \emph{overwritten}---the
continual-learning failure mode studied for mixture-of-experts learners in
\cite{li2024moecl} and in the GAUSE retention analysis \cite{li2026gause}.
The sleeping-experts guarantee says nothing here because the premise ``the
specialist's awake-round losses are fixed targets'' is false; strongly
adaptive algorithms recover the right \emph{weights} quickly but have no
mechanism to ensure there remains a competent expert to weight. \nm{}'s
contribution lives exactly in this gap: the niche memory and the
variance-penalized per-niche objective are mechanisms for making the
``expert exists and is still competent when its regime returns'' premise
\emph{true}, after which adaptive-weighting logic can do its classical job.
\end{remark}

% ============================================================================
\section{Invariant Risk and the Failure of ERM Across Environments}
\label{sec:dd1-invariance}

The last foundation explains \emph{what is being learned inside each niche}.
A niche head sees a handful of episodes of ``its'' regime; those episodes
share the regime's stable structure but differ in incidental, episode-specific
correlations. Pooled empirical risk minimization (ERM) over the episodes is
the default---and it fails in a precise, derivable way: it loads on the
incidental correlations exactly in proportion to how unbalanced the small
sample of episodes happens to be. This section sets up the standard
two-feature construction, derives the failure quantitatively, shows that the
\emph{across-episode variance of per-episode loss} is a statistic that detects
exactly the failure, and positions the \nm{} per-niche objective in the
invariant-risk literature \cite{arjovsky2019irm, krueger2021rex,
peters2016causal}.

\subsection{Environments, Episodes, and Pooled ERM}

\begin{definition}[Multi-environment learning]
\label{def:dd1-environments}
Data arrive grouped into \textbf{episodes} (environments)
$e \in \mathcal{E}$; episode $e$ carries a distribution $P_e$ over
observation--target pairs. A training set consists of samples from $E$
episodes $e_1, \dots, e_E$ drawn from an environment distribution
$\mathcal{Q}$; at deployment a \emph{fresh} episode $e' \sim \mathcal{Q}$ is
encountered. For a predictor $h$ and loss $\ell$, write
$L_e(h) = \E_{P_e}[\ell(h(x), y)]$ for the per-episode risk. \textbf{Pooled
ERM} minimizes $\frac{1}{E}\sum_{j=1}^E L_{e_j}(h)$ (empirically); the
\textbf{invariant} ideal is a predictor whose risk---and whose optimality---is
stable across all $e$ in the support of $\mathcal{Q}$.
\end{definition}

In \nm{}, $\mathcal{E}$ is the set of episodes routed to one niche, $E = E_r
\in \{5, \dots, 7\}$, and $\ell$ is the SPO+ surrogate of
Section~\ref{sec:dd1-spoplus}; for the analysis of this section we use the
squared loss, where everything is exactly computable, and return to the SPO+
case at the end.

\subsection{The Two-Feature Construction}

\begin{definition}[Invariant vs.\ spurious feature model]
\label{def:dd1-twofeature}
Within episode $e$, independent draws are generated as
\begin{equation}
\label{eq:dd1-sem}
x_{\mathrm{inv}} \sim \N(0, 1), \qquad
\varepsilon \sim \N(0, \sigma_\varepsilon^2), \qquad
y = \theta\, x_{\mathrm{inv}} + \varepsilon, \qquad
x_{\mathrm{sp}} = \gamma_e\, y + \nu, \quad \nu \sim \N(0, \sigma_\nu^2),
\end{equation}
with $x_{\mathrm{inv}}, \varepsilon, \nu$ mutually independent. The
\textbf{invariant} coefficient $\theta$ is the same in every episode; the
\textbf{spurious} (anti-causal) coefficient $\gamma_e$ is drawn once per
episode, i.i.d.\ across episodes with $\E[\gamma] = 0$ and
$\Var[\gamma] = \sigma_\gamma^2 > 0$. Write $s^2 := \Var(y) = \theta^2 +
\sigma_\varepsilon^2$. A linear predictor is
$h_{a,b}(x) = a\, x_{\mathrm{inv}} + b\, x_{\mathrm{sp}}$, and the
\textbf{invariant predictor} is $h_{\theta, 0}$.
\end{definition}

\begin{intuition}
$x_{\mathrm{inv}}$ is a genuine driver of $y$ (a stable regime
characteristic); $x_{\mathrm{sp}}$ is a symptom---it is generated \emph{from}
$y$, with an episode-specific gain $\gamma_e$ (an incidental correlation:
e.g.\ a factor that happened to co-move with returns during those particular
weeks). Because $x_{\mathrm{sp}}$ is correlated with $y$ within each episode,
regression happily uses it; because the correlation's sign and size are
episode-specific accidents, doing so buys in-sample fit with out-of-episode
fragility. Averaging the $\gamma_e$'s over infinitely many episodes would
cancel them ($\E\gamma = 0$); the problem is that a niche has $5$--$7$
episodes, and the empirical mean $\bar{\gamma} = \frac{1}{E}\sum_j
\gamma_{e_j}$ is a $\Theta(\sigma_\gamma E^{-1/2})$-sized accident that does
\emph{not} cancel.
\end{intuition}

\begin{proposition}[Pooled-OLS limit loads on the spurious feature]
\label{prop:dd1-pooled-ols}
Consider pooled least squares over $E$ training episodes with equal weights
and infinite within-episode samples (so all second moments are exact). Write
$\bar{\gamma} = \frac{1}{E}\sum_{j} \gamma_{e_j}$,
$\widehat{m}_2 = \frac{1}{E}\sum_j \gamma_{e_j}^2$,
$\widehat{\sigma}_\gamma^2 = \widehat{m}_2 - \bar{\gamma}^2$. The pooled OLS
coefficients $(\widehat{a}, \widehat{b}) = \argmin_{a, b} \frac{1}{E}
\sum_j L_{e_j}(h_{a,b})$ are
\begin{equation}
\label{eq:dd1-pooled-limit}
\widehat{b}
= \frac{\bar{\gamma}\, \sigma_\varepsilon^2}{D},
\qquad
\widehat{a}
= \theta \cdot \frac{\sigma_\nu^2 + s^2\, \widehat{\sigma}_\gamma^2}{D},
\qquad
D := \sigma_\nu^2 + s^2\, \widehat{\sigma}_\gamma^2
+ \bar{\gamma}^2 \sigma_\varepsilon^2 \;>\; 0.
\end{equation}
In particular $\widehat{b} \neq 0$ whenever $\bar{\gamma} \neq 0$: the pooled
solution loads on $x_{\mathrm{sp}}$ \emph{in proportion to the empirical mean
of the episode coefficients}, which for small $E$ is nonzero with
probability one (for continuous $\gamma$) and of typical size
$\sigma_\gamma / \sqrt{E}$.
\end{proposition}

\begin{proof}
Within episode $e$, using \eqref{eq:dd1-sem} and independence:
\[
\E[x_{\mathrm{inv}}^2] = 1, \quad
\E[x_{\mathrm{inv}} x_{\mathrm{sp}}] = \gamma_e \E[x_{\mathrm{inv}} y]
= \gamma_e \theta, \quad
\E[x_{\mathrm{sp}}^2] = \gamma_e^2 s^2 + \sigma_\nu^2, \quad
\E[x_{\mathrm{inv}} y] = \theta, \quad
\E[x_{\mathrm{sp}} y] = \gamma_e s^2 .
\]
Pooling with equal weights averages these moments over the $E$ training
episodes, so the pooled normal equations
$\bar{\Sigma} (\widehat a, \widehat b)^\top = \bar v$ have
\[
\bar{\Sigma} = \begin{pmatrix} 1 & \bar{\gamma}\theta \\
\bar{\gamma}\theta & \widehat{m}_2 s^2 + \sigma_\nu^2 \end{pmatrix},
\qquad
\bar v = \begin{pmatrix} \theta \\ \bar{\gamma} s^2 \end{pmatrix}.
\]
The determinant is
\[
\det \bar{\Sigma}
= \widehat{m}_2 s^2 + \sigma_\nu^2 - \bar{\gamma}^2 \theta^2
= \sigma_\nu^2 + s^2(\widehat{m}_2 - \bar{\gamma}^2)
+ \bar{\gamma}^2 (s^2 - \theta^2)
= \sigma_\nu^2 + s^2 \widehat{\sigma}_\gamma^2
+ \bar{\gamma}^2 \sigma_\varepsilon^2 = D > 0 .
\]
By Cramer's rule,
\[
\widehat{b} = \frac{1 \cdot \bar{\gamma} s^2 - \bar{\gamma}\theta \cdot
\theta}{D} = \frac{\bar{\gamma}(s^2 - \theta^2)}{D}
= \frac{\bar{\gamma}\, \sigma_\varepsilon^2}{D},
\qquad
\widehat{a} = \frac{(\widehat{m}_2 s^2 + \sigma_\nu^2)\theta -
\bar{\gamma}\theta \cdot \bar{\gamma} s^2}{D}
= \frac{\theta\bigl( \sigma_\nu^2 + s^2 \widehat{\sigma}_\gamma^2
\bigr)}{D}. \qedhere
\]
\end{proof}

Note the mechanism: $\widehat{b} \propto \bar{\gamma}$, an average of
mean-zero accidents that a law of large numbers would kill---\emph{if $E$
were large}. ERM is not wrong asymptotically; it is wrong at $E = 5$--$7$,
which is the only regime \nm{} lives in. Pooled ERM also genuinely
\emph{prefers} this loading in-sample:

\begin{lemma}[ERM strictly descends into the spurious direction]
\label{lem:dd1-erm-descends}
The per-episode risk of $h_{a,b}$ is
\begin{equation}
\label{eq:dd1-Le}
L_e(h_{a,b})
= \bigl( \theta - a - b\gamma_e\theta \bigr)^2
+ \bigl( 1 - b\gamma_e \bigr)^2 \sigma_\varepsilon^2
+ b^2 \sigma_\nu^2,
\end{equation}
and at the invariant predictor $(a, b) = (\theta, 0)$ the pooled risk has
$\frac{\partial}{\partial b}\, \frac{1}{E}\sum_j L_{e_j} =
-2\bar{\gamma}\sigma_\varepsilon^2$. Hence whenever $\bar{\gamma} \neq 0$,
moving $b$ slightly in the direction of $\operatorname{sign}(\bar\gamma)$
strictly reduces pooled training risk below the invariant predictor's.
\end{lemma}

\begin{proof}
The residual is $y - a x_{\mathrm{inv}} - b x_{\mathrm{sp}}
= (\theta - a - b\gamma_e \theta) x_{\mathrm{inv}} + (1 - b\gamma_e)
\varepsilon - b \nu$ (substitute $x_{\mathrm{sp}} = \gamma_e(\theta
x_{\mathrm{inv}} + \varepsilon) + \nu$); squaring and using independence and
unit/known variances gives \eqref{eq:dd1-Le}. Differentiate in $b$:
\[
\frac{\partial L_e}{\partial b}
= -2\gamma_e\theta\,(\theta - a - b\gamma_e\theta)
- 2\gamma_e \sigma_\varepsilon^2 (1 - b\gamma_e) + 2 b \sigma_\nu^2
\;\xrightarrow{\;(a,b) = (\theta, 0)\;}\;
0 - 2\gamma_e \sigma_\varepsilon^2 + 0 .
\]
Averaging over the training episodes gives $-2\bar\gamma
\sigma_\varepsilon^2$, nonzero iff $\bar\gamma \neq 0$.
\end{proof}

\begin{proposition}[Expected excess risk on a fresh episode]
\label{prop:dd1-excess}
For any fixed predictor $h_{a,b}$ and a fresh episode $e' \sim \mathcal{Q}$
(i.e.\ $\gamma_{e'}$ a fresh draw with mean $0$ and variance
$\sigma_\gamma^2$, independent of the training episodes),
\begin{equation}
\label{eq:dd1-excess}
\E_{\gamma_{e'}}\bigl[ L_{e'}(h_{a,b}) \bigr] - L_{e'}(h_{\theta, 0})
= (\theta - a)^2 \;+\; b^2 \bigl( \sigma_\nu^2 + s^2 \sigma_\gamma^2 \bigr)
\;\ge\; 0,
\end{equation}
with equality iff $(a,b) = (\theta, 0)$; note $L_{e'}(h_{\theta,0}) =
\sigma_\varepsilon^2$ deterministically. For the pooled-OLS solution
\eqref{eq:dd1-pooled-limit}, the leading term for small $|\bar\gamma|$ is
\begin{equation}
\label{eq:dd1-excess-rate}
\widehat{b}^{\,2}\bigl( \sigma_\nu^2 + s^2\sigma_\gamma^2 \bigr)
= \frac{\bar{\gamma}^2\, \sigma_\varepsilon^4 \,(\sigma_\nu^2 +
s^2\sigma_\gamma^2)}{D^2},
\qquad
\E\bigl[\bar{\gamma}^2\bigr] = \frac{\sigma_\gamma^2}{E},
\end{equation}
while $(\theta - \widehat a)^2 = O(\bar\gamma^4)$. The expected damage of
pooling thus decays only as $1/E$---negligible for hundreds of episodes,
material for $E = 5$--$7$.
\end{proposition}

\begin{proof}
Take the expectation of \eqref{eq:dd1-Le} over $\gamma_{e'}$ with
$\E\gamma_{e'} = 0$, $\E\gamma_{e'}^2 = \sigma_\gamma^2$:
\[
\E\bigl[(\theta - a - b\gamma_{e'}\theta)^2\bigr]
= (\theta - a)^2 + b^2\theta^2\sigma_\gamma^2,
\qquad
\E\bigl[(1 - b\gamma_{e'})^2\bigr]\sigma_\varepsilon^2
= \sigma_\varepsilon^2 + b^2 \sigma_\varepsilon^2 \sigma_\gamma^2 .
\]
Summing with $b^2\sigma_\nu^2$ and subtracting
$L_{e'}(h_{\theta,0}) = \sigma_\varepsilon^2$ (immediate from
\eqref{eq:dd1-Le} at $(\theta, 0)$) gives \eqref{eq:dd1-excess}, using
$\theta^2 + \sigma_\varepsilon^2 = s^2$. For the rate: $\widehat b =
\bar\gamma \sigma_\varepsilon^2 / D$ from
Proposition~\ref{prop:dd1-pooled-ols}, and $\theta - \widehat a =
\theta \bar\gamma^2 \sigma_\varepsilon^2 / D$ (subtract
\eqref{eq:dd1-pooled-limit} from $\theta$), whose square is $O(\bar\gamma^4)$.
Finally $\E[\bar\gamma^2] = \Var(\bar\gamma) = \sigma_\gamma^2 / E$ for
i.i.d.\ mean-zero $\gamma_{e_j}$.
\end{proof}

\begin{example}[Numbers at niche scale]
\label{ex:dd1-pooled-numbers}
Take $\theta = 1$, $\sigma_\varepsilon^2 = 0.5$ (so $s^2 = 1.5$),
$\sigma_\nu^2 = 1$, and $E = 6$ training episodes with
$\gamma = (0.4, -0.2, 0.3, -0.1, 0.2, 0.0)$---a sample from a mean-zero law
that happens, as small samples do, to lean positive: $\bar\gamma = 0.1$,
$\widehat{m}_2 = 0.34/6 \approx 0.0567$, $\widehat\sigma_\gamma^2 \approx
0.0467$. Then $D = 1 + 1.5 \times 0.0467 + 0.01 \times 0.5 \approx 1.075$ and
\[
\widehat{b} = \frac{0.1 \times 0.5}{1.075} \approx 0.047, \qquad
\widehat{a} \approx 1 \times \frac{1.070}{1.075} \approx 0.995 .
\]
With true $\sigma_\gamma^2 = 0.05$, the expected fresh-episode excess
\eqref{eq:dd1-excess} is approximately $0.047^2 (1 + 1.5 \times 0.05) +
0.005^2 \approx 0.0023$---about $0.5\%$ of the irreducible risk
$\sigma_\varepsilon^2 = 0.5$ per fresh episode, incurred systematically, with
the sign of the error flipping adversarially on episodes where
$\gamma_{e'}$ opposes $\bar\gamma$.
\end{example}

\subsection{The Variance Statistic Detects Exactly the Loading}

The decisive structural fact about the construction is visible in
\eqref{eq:dd1-Le}: the per-episode risk depends on the episode \emph{only
through the products $b\gamma_e$}. A predictor that ignores the spurious
feature has the \emph{same} risk in every episode; a predictor that uses it
has episode-dependent risk. Across-episode variance is therefore a
\emph{detector} of spurious loading.

\begin{proposition}[Variance identifies invariance]
\label{prop:dd1-variance-detects}
Let $\Var_e[\,\cdot\,]$ denote the empirical variance across the training
episodes $e_1, \dots, e_E$, and suppose at least two of the
$\gamma_{e_j}$ are distinct. Then for the model
\eqref{eq:dd1-sem}:
\begin{enumerate}
\item If $b = 0$, then $L_{e_j}(h_{a, 0}) = (\theta - a)^2 +
\sigma_\varepsilon^2$ for every $j$, so $\Var_e\bigl[ L_e(h_{a,0}) \bigr] =
0$.
\item If $b \neq 0$, then $\Var_e\bigl[ L_e(h_{a,b}) \bigr] > 0$, except for
at most finitely many exceptional values of $(a, b)$ (those making the
quadratic $\gamma \mapsto L(\gamma)$ in \eqref{eq:dd1-Le} take equal values
at all the distinct observed $\gamma_{e_j}$, possible only when fewer than
three distinct values are observed and $(a,b)$ is tuned to the pair).
With three or more distinct $\gamma_{e_j}$, $b \neq 0$ implies
$\Var_e[L_e] > 0$ without exception.
\end{enumerate}
Consequently the population of zero-variance predictors is exactly the
$b = 0$ axis (generically), on which the pooled mean risk is minimized at
the invariant predictor $a = \theta$: minimizing
$\Var_e[L_e] + \beta\, \E_e[L_e]$ with $\beta > 0$ and a sufficiently large
variance weight recovers $h_{\theta, 0}$, whereas pooled ERM
($\beta$-only) strictly prefers $\widehat b \neq 0$ by
Lemma~\ref{lem:dd1-erm-descends}.
\end{proposition}

\begin{proof}
(1) Set $b = 0$ in \eqref{eq:dd1-Le}: every $\gamma_e$-dependent term carries
a factor $b$, leaving $(\theta - a)^2 + \sigma_\varepsilon^2$, constant
across episodes; the empirical variance of a constant list is $0$.

(2) Regard \eqref{eq:dd1-Le} as a polynomial in $\gamma$:
\[
L(\gamma) = \underbrace{b^2\theta^2 \,\gamma^2 - 2b\theta(\theta - a)\gamma
+ (\theta - a)^2}_{\text{first term}}
+ \underbrace{\sigma_\varepsilon^2 b^2 \gamma^2 - 2\sigma_\varepsilon^2 b
\gamma + \sigma_\varepsilon^2}_{\text{second term}} + b^2\sigma_\nu^2,
\]
with leading coefficient $b^2(\theta^2 + \sigma_\varepsilon^2) = b^2 s^2 > 0$
when $b \neq 0$: a nondegenerate quadratic. The empirical variance vanishes
iff $L(\gamma_{e_j})$ is the same number for all $j$. A nondegenerate
quadratic takes any given value at most twice, so if three or more distinct
$\gamma_{e_j}$ are observed, equality of all values is impossible and
$\Var_e[L_e] > 0$. With exactly two distinct values $\gamma \neq \gamma'$,
equality forces the quadratic's axis of symmetry to bisect them, i.e.\ pins
one algebraic relation between $a$ and $b$---the finitely-many-exceptions
caveat. The final claims restate Lemma~\ref{lem:dd1-erm-descends} and part
(1) together with the strict positivity in part (2): for large variance
weight, any $b \neq 0$ is dominated by its $b = 0$ projection, and among
$b = 0$ predictors the mean term selects $a = \theta$.
\end{proof}

\begin{keypoint}
The variance statistic is not a generic regularizer; in this construction it
is a \emph{consistent test} for the precise pathology of
Proposition~\ref{prop:dd1-pooled-ols}. Pooled ERM cannot see the difference
between fitting structure and fitting accidents, because both reduce mean
training risk (Lemma~\ref{lem:dd1-erm-descends}); the across-episode variance
sees \emph{only} the accidents, because structure---by definition of being
shared across the niche's episodes---contributes nothing to it
(Proposition~\ref{prop:dd1-variance-detects}). This is also exactly the
quantity the small-sample transfer story of
Section~\ref{sec:dd1-concentration} wants minimized: a smaller
across-episode spread is what would tighten an empirical-Bernstein-style
certificate (Remark~\ref{rem:dd1-small-n}) and what shrinks the variance in
Cantelli's one-sided bound (Theorem~\ref{thm:dd1-cantelli}).
\end{keypoint}

\subsection{Relation to IRM, REx, and Causal Invariance; the \nm{} Objective}

Three strands of the invariance literature meet here.
\textbf{Invariant causal prediction} \cite{peters2016causal} formalizes the
target: under interventions that vary across environments, the conditional
distribution of the response given its \emph{causal parents} is invariant,
and (under conditions) the causal feature set is identifiable as the one
producing invariant residuals. In Definition~\ref{def:dd1-twofeature},
$x_{\mathrm{inv}}$ is a parent and $x_{\mathrm{sp}}$ a child of $y$; the
invariant predictor is the causal one.
\textbf{IRM} \cite{arjovsky2019irm} pursues this target through a bilevel
program---learn a representation such that one classifier on top of it is
simultaneously per-environment optimal---in practice relaxed to a gradient
penalty $\sum_e \|\nabla_{w \mid w = 1} L_e(w \cdot h)\|^2$.
\textbf{REx} \cite{krueger2021rex} observes that a simpler scalarization
often suffices and is better behaved: penalize the spread of per-environment
risks directly, minimizing
$\mathrm{Var}_e[L_e(h)] + \beta\, \E_e[L_e(h)]$-type objectives (V-REx);
robustness to the spread of risks extrapolates to environments whose
``accidents'' are larger than any seen in training.

The \nm{} per-niche objective is the V-REx scalarization applied to a
\emph{decision-focused} episode risk. For niche $r$ with stored episodes
$\mathcal{E}_r$ and head parameters $\theta_r$, define the SPO+ episode risk
$R_e(\theta_r) = \frac{1}{|e|}\sum_{(x, y) \in e}
\ell_{\mathrm{SPO+}}\bigl( \hat y_{\theta_r}(x);\, y \bigr)$ and train by
\begin{equation}
\label{eq:dd1-nm-objective}
J_r(\theta_r)
\;=\; \Var_{e \in \mathcal{E}_r}\bigl[ R_e(\theta_r) \bigr]
\;+\; \beta\; \E_{e \in \mathcal{E}_r}\bigl[ R_e(\theta_r) \bigr],
\qquad \beta > 0 .
\end{equation}
Two substitutions distinguish \eqref{eq:dd1-nm-objective} from textbook REx.
First, the risk inside the variance is the SPO+ \emph{decision} risk, so
``invariance'' is demanded of the quantity that matters
(Section~\ref{sec:dd1-dfl}): a head may have episode-varying \emph{forecast}
error and still be exactly invariant in decision quality, and the objective
correctly declines to punish it. Second, the environments are not given
exogenously but are the episodes routed to the niche---the regime-indexed
environment structure of the invariant predict-and-optimize formulation
\cite{yan2026invpnco}---so the variance is computed within a (putative)
regime, where invariance is actually achievable, rather than across all of
history, where it is not.

\begin{warning}
\textbf{Variance alone ($\beta = 0$) is degenerate.} The zero-variance set
contains far more than invariant-and-good predictors: \emph{any} head whose
episode risks are constant achieves $J_r = 0$ at $\beta = 0$, including
uniformly terrible ones. Concretely, a head whose forecasts induce the same
fixed bad portfolio in all states has $R_e \equiv$ (the same poor value) on
every episode---zero variance, maximal complacency; in the squared-loss model
of Proposition~\ref{prop:dd1-variance-detects}, the entire plane $b = 0$ is
zero-variance regardless of how wrong $a$ is. The mean term is therefore not
a tuning nicety but a load-bearing component: it is the only part of
\eqref{eq:dd1-nm-objective} that distinguishes invariant-good from
invariant-bad. This degeneracy predicts, qualitatively and in advance, the
$\beta = 0$ collapse observed in the companion experiments: post-reactivation
decision regret $0.0356$ for the variance-only objective versus $0.0321$ for
plain retained ERM---the ablation does not merely lose the invariance
benefit, it underperforms doing nothing clever at all, exactly as a
degenerate objective should.
\end{warning}

\begin{remark}[What carries over from squared loss to SPO+]
\label{rem:dd1-carryover}
The exact algebra of
Propositions~\ref{prop:dd1-pooled-ols}--\ref{prop:dd1-variance-detects} is
special to squared loss, where risks are quadratics in $\gamma_e$. What
transfers to the SPO+ risk is the structural trichotomy, none of whose
ingredients used quadraticity: (i) pooled ERM's first-order preference for
episode-specific signal scales with the \emph{empirical mean} of the
accident process and is therefore $\Theta(E^{-1/2})$-sized at niche scale;
(ii) any loading on episode-specific signal makes the per-episode risk
\emph{episode-dependent}, hence visible to the across-episode variance;
(iii) zero variance plus low mean jointly---and only jointly---certify
``invariantly good.'' These three facts, with SPO+ risks in place of squared
risks and with the boundedness $R_e$-analogues supplied by
Proposition~\ref{prop:dd1-regret-range}, are the form in which Part II uses
this section.
\end{remark}

This completes the toolkit. Part II assembles it: the decision layer and
SPO+ machinery of Section~\ref{sec:dd1-dfl} define what each niche head
optimizes \eqref{eq:dd1-nm-objective}; the EG gate of
Section~\ref{sec:dd1-online} allocates episodes and training resources to
niches, deliberately forgoing the adversarial guarantee
(Section~\ref{sec:dd1-online-limits}) in favor of analyzable replicator
dynamics; and the retention and transfer theorems combine the concentration
inequalities of Section~\ref{sec:dd1-concentration} with the Pinsker bridge
of Section~\ref{sec:dd1-distances} to state precisely when a dormant niche's
certified competence survives until its regime returns---the capability gap,
identified in Remark~\ref{rem:dd1-gap}, that the sleeping-experts literature
leaves open and \nm{} is built to close.
